% !TEX root = 第六章_定积分(答案版).tex
\setcounter{chapter}{5}
\pagestyle{fancy}

\chapter{定积分}


\section{积分计算}


\subsection{对称性}

\begin{proposition}[\marktwostar]
    记
    \[
        I(m,n)=\int_0^{\frac{\pi}{2}}\cos^m x\sin^n x\,\mathrm{d}x,
    \]
    则 $I(m,n)=I(n,m)$,且
    \[
        I(m,n)=\dfrac{m-1}{m+n}I(m-2,n)
        =\dfrac{n-1}{m+n}I(m,n-2),
    \]
    并有
    \[
        I(m,n)=
        \begin{cases}
            \dfrac{(m-1)!!(n-1)!!}{(m+n)!!},               & m,n\text{ 不全为偶数}, \\[6pt]
            \dfrac{(m-1)!!(n-1)!!}{(m+n)!!}\dfrac{\pi}{2}, & m,n\text{ 都为偶数}.
        \end{cases}
    \]
\end{proposition}

\begin{theorem}[\markstar]
    设 $f(x),g(x)$ 为两个实函数,那么:
    \begin{enumerate}
        \item $f(x)$ 与 $g(x)$ 都关于 $x=x_0$ 轴对称,则 $f(x)g(x)$ 也关于 $x=x_0$ 轴对称;
        \item $f(x)$ 关于 $x=x_0$ 轴对称,$g(x)$ 关于 $(x_0,0)$ 中心对称,则 $f(x)g(x)$ 关于 $(x_0,0)$ 中心对称;
        \item $f(x)$ 与 $g(x)$ 都关于 $(x_0,0)$ 中心对称,则 $f(x)g(x)$ 关于 $x=x_0$ 轴对称.
    \end{enumerate}
\end{theorem}

\begin{corollary}[\markstar]
    对任意非负整数 $m,n$,有:
    \begin{enumerate}
        \item $\sin^n x$ 关于 $x=\dfrac{\pi}{2}$ 轴对称;
        \item $m$ 为偶数时,$\cos^m x$ 关于 $x=\dfrac{\pi}{2}$ 轴对称;$m$ 为奇数时,$\cos^m x$ 关于 $\left(\dfrac{\pi}{2},0\right)$ 中心对称;
        \item $m$ 为偶数时,$\cos^m x\sin^n x$ 关于 $x=\dfrac{\pi}{2}$ 轴对称;当 $m$ 为奇数时,$\cos^m x\sin^n x$ 关于 $\left(\dfrac{\pi}{2},0\right)$ 中心对称.
    \end{enumerate}
\end{corollary}

\begin{proposition}[\markstar]
    对于任意非负整数 $m,n$,有
    \[
        J(m,n)\overset{\triangle}{=}
        \int_0^\pi\cos^m x\sin^n x\,\mathrm{d}x
        =\begin{cases}
            0,       & m\text{ 为奇数}, \\
            2I(m,n), & m\text{ 为偶数}.
        \end{cases}
    \]
\end{proposition}

\begin{remark}
    此时关于 $m,n$ 并非对称,即一般情况下 $J(m,n)\neq J(n,m)$.
\end{remark}

\begin{proposition}[\markstar]
    对于任意非负整数 $m,n$,有
    \[
        K(m,n)\overset{\triangle}{=}
        \int_0^{2\pi}\cos^m x\sin^n x\,\mathrm{d}x
        =\begin{cases}
            0,       & n\text{ 为奇数}, \\
            2J(m,n), & n\text{ 为偶数},
        \end{cases}
    \]
    从而
    \[
        K(m,n)=\begin{cases}
            0,       & m,n\text{ 不全为偶数}, \\
            4I(m,n), & m,n\text{ 都为偶数}.
        \end{cases}
    \]
\end{proposition}

%例题6.1.1
\begin{example}
    计算积分:
    \begin{enumerate}
        \item $\displaystyle\int_0^{\frac{\pi}{2}}\cos^3x\sin^2x\,\mathrm{d}x$;
        \item $\displaystyle\int_0^{\frac{\pi}{2}}\cos^4x\sin^3x\,\mathrm{d}x$;
        \item $\displaystyle\int_0^{\frac{\pi}{2}}\sin^3x\cos^5x\,\mathrm{d}x$;
        \item $\displaystyle\int_0^{\frac{\pi}{2}}\sin^2x\cos^4x\,\mathrm{d}x$.
    \end{enumerate}
\end{example}
\exampleblank

\begin{solution}
    使用公式
    \[
        \int_0^{\frac{\pi}{2}}\sin^m x\cos^n x\,\mathrm{d}x
        =\dfrac{1}{2}B\left(\dfrac{m+1}{2},\dfrac{n+1}{2}\right).
    \]
    逐项代入得到
    \[
        \boxed{\dfrac{2}{15}},\qquad
        \boxed{\dfrac{2}{35}},\qquad
        \boxed{\dfrac{1}{24}},\qquad
        \boxed{\dfrac{\pi}{32}}.
    \]
\end{solution}

%例题6.1.2
\begin{example}
    计算积分:
    \begin{enumerate}
        \item $\displaystyle\int_0^\pi\cos^3x\sin^2x\,\mathrm{d}x$;
        \item $\displaystyle\int_0^\pi\cos^4x\sin^3x\,\mathrm{d}x$;
        \item $\displaystyle\int_0^\pi\cos^n x\,\mathrm{d}x$;
        \item $\displaystyle\int_0^\pi\sin^n x\,\mathrm{d}x$.
    \end{enumerate}
\end{example}
\exampleblank

\begin{solution}
    \begin{enumerate}
        \item 令 $x\mapsto\pi-x$,被积函数变号,故积分为 $0$.
        \item 关于 $\dfrac{\pi}{2}$ 对称后,
              \[
                  \int_0^\pi\cos^4x\sin^3x\,\mathrm{d}x
                  =2\int_0^{\frac{\pi}{2}}\cos^4x\sin^3x\,\mathrm{d}x
                  =\boxed{\dfrac{4}{35}}.
              \]
        \item 当 $n$ 为奇数时,由关于 $\dfrac{\pi}{2}$ 的反对称性,积分为 $0$.当 $n=2m$ 时,
              \[
                  \int_0^\pi\cos^{2m}x\,\mathrm{d}x
                  =\boxed{\dfrac{\pi C_m^{2m}}{4^m}}.
              \]
        \item 由对称性及 Wallis 公式,
              \[
                  \int_0^\pi\sin^n x\,\mathrm{d}x
                  =
                  \begin{cases}
                      \dfrac{\pi C_m^{2m}}{4^m},       & n=2m,   \\[6pt]
                      \dfrac{2^{2m+1}(m!)^2}{(2m+1)!}, & n=2m+1.
                  \end{cases}
              \]
    \end{enumerate}
\end{solution}

%例题6.1.3
\begin{example}
    计算积分:
    \begin{enumerate}
        \item $\displaystyle\int_0^{2\pi}\cos^4x\sin^3x\,\mathrm{d}x$;
        \item $\displaystyle\int_0^{2\pi}\sin^3x\cos^5x\,\mathrm{d}x$;
        \item $\displaystyle\int_0^{2\pi}\sin^2x\cos^4x\,\mathrm{d}x$;
        \item $\displaystyle\int_0^{2\pi}\sin^n x\,\mathrm{d}x
                  =\int_0^{2\pi}\cos^n x\,\mathrm{d}x$.
    \end{enumerate}
\end{example}
\exampleblank

\begin{solution}
    \begin{enumerate}
        \item 被积函数在一个整周期上的积分为零,故结果为 $\boxed{0}$.
        \item 在 $\lbrack0,\pi\rbrack$ 上作代换 $x\mapsto\pi-x$,被积函数变号,故一个周期上的积分为 $\boxed{0}$.
        \item 四个象限贡献相同,因此
              \[
                  \int_0^{2\pi}\sin^2x\cos^4x\,\mathrm{d}x
                  =4\int_0^{\frac{\pi}{2}}\sin^2x\cos^4x\,\mathrm{d}x
                  =\boxed{\dfrac{\pi}{8}}.
              \]
        \item 两积分由平移 $x\mapsto x+\dfrac{\pi}{2}$ 相等.若 $n$ 为正整数,则公共值为
              \[
                  \begin{cases}
                      0,                          & n\text{ 为奇数}, \\[4pt]
                      \dfrac{2\pi C_m^{2m}}{4^m}, & n=2m.
                  \end{cases}
              \]
    \end{enumerate}
\end{solution}

%例题6.1.4
\begin{example}
    计算积分:
    \begin{enumerate}
        \item $\displaystyle\int_{-\pi}^{\pi}\cos^4x\sin^4x\,\mathrm{d}x$;
        \item $\displaystyle\int_{-\pi}^{3\pi}\sin^5x\cos^6x\,\mathrm{d}x$;
        \item $\displaystyle\int_{-\pi}^{0}\cos^3x\sin^2x\,\mathrm{d}x$;
        \item $\displaystyle\int_0^\pi|\cos^5x|\,\mathrm{d}x$.
    \end{enumerate}
\end{example}
\exampleblank

\begin{solution}
    \begin{enumerate}
        \item 由周期性及
              \[
                  \sin^4x\cos^4x=\dfrac{1}{16}\sin^4(2x),
              \]
              得
              \[
                  \int_{-\pi}^{\pi}\cos^4x\sin^4x\,\mathrm{d}x
                  =\boxed{\dfrac{3\pi}{64}}.
              \]
        \item 被积函数在平移 $\pi$ 后变号,积分区间包含两个完整的 $2\pi$ 周期,故结果为 $\boxed{0}$.
        \item 由偶性把积分化为 $\displaystyle\int_0^\pi\cos^3x\sin^2x\,\mathrm{d}x$,再关于 $\dfrac{\pi}{2}$ 配对,结果为 $\boxed{0}$.
        \item 关于 $\dfrac{\pi}{2}$ 对称,故
              \[
                  \int_0^\pi|\cos^5x|\,\mathrm{d}x
                  =2\int_0^{\frac{\pi}{2}}\cos^5x\,\mathrm{d}x
                  =\boxed{\dfrac{16}{15}}.
              \]
    \end{enumerate}
\end{solution}

%例题6.1.5
\begin{example}
    计算积分:
    \begin{enumerate}
        \item $\displaystyle\int_0^\pi\dfrac{\mathrm{d}x}{1+\sin^2x}$;
        \item $\displaystyle\int_0^{2\pi}\dfrac{\mathrm{d}x}{1+a\cos x}\quad(-1<a<1)$.
    \end{enumerate}
\end{example}
\exampleblank

\begin{solution}
    \begin{enumerate}
        \item 关于 $\dfrac{\pi}{2}$ 对称并令 $t=\tan x$,得到
              \[
                  \begin{aligned}
                      \int_0^\pi\dfrac{\mathrm{d}x}{1+\sin^2x}
                       & =2\int_0^{\frac{\pi}{2}}\dfrac{\mathrm{d}x}{1+\sin^2x} \\
                       & =2\int_0^{+\infty}\dfrac{\mathrm{d}t}{1+2t^2}
                      =\boxed{\dfrac{\pi}{\sqrt2}}.
                  \end{aligned}
              \]
        \item 使用 $t=\tan\dfrac{x}{2}$,并把 $0$ 到 $2\pi$ 分成两段,可得标准积分
              \[
                  \int_0^{2\pi}\dfrac{\mathrm{d}x}{1+a\cos x}
                  =\boxed{\dfrac{2\pi}{\sqrt{1-a^2}}},
                  \qquad |a|<1.
              \]
    \end{enumerate}
\end{solution}

%例题6.1.6
\begin{example}
    利用对称性,计算下列积分:
    \begin{enumerate}
        \item $\displaystyle\int_0^\pi\dfrac{x}{1+\cos^2x}\,\mathrm{d}x$;
        \item $\displaystyle\int_0^1\dfrac{x}{e^x+e^{1-x}}\,\mathrm{d}x$;
        \item $\displaystyle\int_{-2}^{2}x\ln(1+e^x)\,\mathrm{d}x$;
        \item $\displaystyle\int_0^{\frac{\pi}{4}}\ln(1+\tan x)\,\mathrm{d}x$;
        \item $\displaystyle\int_0^{\frac{\pi}{2}}\dfrac{\sin^n x}{\sin^n x+\cos^n x}\,\mathrm{d}x$;
        \item $\displaystyle\int_0^\pi
                  \dfrac{a^n\sin^2x+b^n\cos^2x}{a^{2n}\sin^2x+b^{2n}\cos^2x}\,\mathrm{d}x$.
    \end{enumerate}
\end{example}
\exampleblank

\begin{solution}
    \begin{enumerate}
        \item 令 $x\mapsto\pi-x$ 并与原式相加:
              \[
                  2I=\pi\int_0^\pi\dfrac{\mathrm{d}x}{1+\cos^2x}
                  =\dfrac{\pi^2}{\sqrt2}.
              \]
              故 $I=\boxed{\dfrac{\pi^2}{2\sqrt2}}$.

        \item 记被积函数分母为 $D(x)=e^x+e^{1-x}$.因 $D(1-x)=D(x)$,
              \[
                  2I=\int_0^1\dfrac{\mathrm{d}x}{D(x)}
                  =e^{-\frac{1}{2}}\int_0^{\frac{1}{2}}\operatorname{sech}u\,\mathrm{d}u.
              \]
              因此
              \[
                  I=\boxed{\dfrac{e^{-\frac{1}{2}}}{2}
                      \arctan\left(\sinh\dfrac{1}{2}\right)}.
              \]

        \item 将正,负半轴配对,并使用
              \[
                  \ln(1+e^x)-\ln(1+e^{-x})=x,
              \]
              得
              \[
                  I=\int_0^2x^2\,\mathrm{d}x=\boxed{\dfrac{8}{3}}.
              \]

        \item 令 $x\mapsto\dfrac{\pi}{4}-x$.由于
              \[
                  \ln(1+\tan x)
                  +\ln\left(1+\tan\left(\dfrac{\pi}{4}-x\right)\right)=\ln2,
              \]
              故
              \[
                  I=\boxed{\dfrac{\pi}{8}\ln2}.
              \]

        \item 令 $x\mapsto\dfrac{\pi}{2}-x$,所得两个被积函数之和为 $1$,所以
              \[
                  I=\boxed{\dfrac{\pi}{4}}.
              \]

        \item 记 $A=a^n$,$B=b^n$.被积函数以 $\pi$ 为周期,且令 $t=\tan x$ 后
              \[
                  \begin{aligned}
                      \int_0^{\frac{\pi}{2}}
                      \dfrac{A\sin^2x+B\cos^2x}
                      {A^2\sin^2x+B^2\cos^2x}\,\mathrm{d}x
                       & =\int_0^{+\infty}
                      \dfrac{At^2+B}{(A^2t^2+B^2)(1+t^2)}\,\mathrm{d}t \\
                       & =\dfrac{\pi}{A+B}.
                  \end{aligned}
              \]
              因而原积分为
              \[
                  \boxed{\dfrac{2\pi}{a^n+b^n}}.
              \]
    \end{enumerate}
\end{solution}

%例题6.1.7
\begin{example}
    设 $f\in C(\lbrack 0,a\rbrack)$,$a>0$.
    \begin{enumerate}
        \item 在 $\lbrack 0,a\rbrack$ 上 $f(x)+f(a-x)\neq0$,计算
              \[
                  I=\int_0^a\dfrac{f(x)}{f(x)+f(a-x)}\,\mathrm{d}x;
              \]
        \item 在 $\lbrack 0,a\rbrack$ 上 $f(x)f(a-x)\equiv1$,计算
              \[
                  I=\int_0^a\dfrac{\mathrm{d}x}{1+f(x)}.
              \]
    \end{enumerate}
\end{example}
\exampleblank

\begin{solution}
    \begin{enumerate}
        \item 作代换 $x\mapsto a-x$,得到
              \[
                  I=\int_0^a\dfrac{f(a-x)}{f(x)+f(a-x)}\,\mathrm{d}x.
              \]
              与原式相加得 $2I=a$,故 $I=\boxed{\dfrac{a}{2}}$.

        \item 同样作代换 $x\mapsto a-x$,再用 $f(x)f(a-x)=1$,得到
              \[
                  I=\int_0^a\dfrac{f(x)}{1+f(x)}\,\mathrm{d}x.
              \]
              与原式相加仍有 $2I=a$,故 $I=\boxed{\dfrac{a}{2}}$.
    \end{enumerate}
\end{solution}

%例题6.1.8
\begin{example}
    设 $f$ 为连续函数,证明下列等式:
    \begin{enumerate}
        \item
              \[
                  \int_1^a f\left(x^2+\dfrac{a^2}{x^2}\right)\dfrac{\mathrm{d}x}{x}
                  =\int_1^a f\left(x+\dfrac{a^2}{x}\right)\dfrac{\mathrm{d}x}{x};
              \]
        \item
              \[
                  \int_0^{2\pi}f(a\cos x+b\sin x)\,\mathrm{d}x
                  =2\int_0^\pi f\left(\sqrt{a^2+b^2}\cos x\right)\,\mathrm{d}x.
              \]
    \end{enumerate}
\end{example}
\exampleblank

\begin{solution}
    \begin{enumerate}
        \item 左端令 $t=x^2$,得
              \[
                  \dfrac{1}{2}\int_1^{a^2}
                  f\left(t+\dfrac{a^2}{t}\right)\dfrac{\mathrm{d}t}{t}.
              \]
              将积分在 $t=a$ 处分开,并在第二段作代换 $t=\dfrac{a^2}{u}$,两段相等,故结果为
              \[
                  \int_1^a f\left(u+\dfrac{a^2}{u}\right)\dfrac{\mathrm{d}u}{u}.
              \]

        \item 取 $R=\sqrt{a^2+b^2}$ 及角 $\varphi$,使
              \[
                  a\cos x+b\sin x=R\cos(x-\varphi).
              \]
              利用 $2\pi$ 周期性平移积分区间,左端等于
              \[
                  \int_0^{2\pi}f(R\cos x)\,\mathrm{d}x.
              \]
              再把区间分成两个长度为 $\pi$ 的部分,并在后一部分作 $x\mapsto2\pi-x$,得到题设右端.
    \end{enumerate}
\end{solution}



\subsection{积分与极限}

%例题6.1.9
\begin{example}
    设 $f(x)$ 在 $\lbrack 0,1\rbrack$ 上连续,求证:
    \[
        \lim\limits_{n\to\infty}\dfrac{1}{n}
        \sum\limits_{k=1}^{n-1}(-1)^{k+1}f\left(\dfrac{k}{n}\right)=0.
    \]
\end{example}
\exampleblank

\begin{solution}
    设 $\omega_f(\delta)$ 为 $f$ 在 $\lbrack0,1\rbrack$ 上的连续模.将相邻两项配对.若 $n=2m$,则
    \[
        \begin{aligned}
            \left|\dfrac{1}{n}\sum\limits_{k=1}^{n-1}
            (-1)^{k+1}f\left(\dfrac{k}{n}\right)\right|
             & \leq\dfrac{1}{n}\sum\limits_{j=1}^{m-1}
            \left|f\left(\dfrac{2j-1}{n}\right)
            -f\left(\dfrac{2j}{n}\right)\right|
            +\dfrac{\|f\|_\infty}{n}                             \\
             & \leq\dfrac{1}{2}\omega_f\left(\dfrac{1}{n}\right)
            +\dfrac{\|f\|_\infty}{n}.
        \end{aligned}
    \]
    奇数情形只多出一个同样趋于零的端项.由一致连续性,右端趋于零,故原极限为 $0$.
\end{solution}

%例题6.1.10
\begin{example}
    求极限
    \[
        \lim\limits_{n\to\infty}
        \left(b^{\frac{1}{n}}-1\right)
        \sum\limits_{i=0}^{n-1}
        b^{\frac{1}{n}}\sin b^{\frac{1}{n}}
        \sin b^{\frac{n+1}{2n}}.
    \]
\end{example}
\exampleblank

\begin{solution}
    按当前题面,求和号内的表达式不含指标 $i$,因此该和恰为其单项的 $n$ 倍.若 $b>0$,则
    \[
        n\left(b^{\frac{1}{n}}-1\right)\longrightarrow\ln b,
        \qquad
        b^{\frac{1}{n}}\longrightarrow1,
        \qquad
        b^{\frac{n+1}{2n}}\longrightarrow\sqrt b.
    \]
    所以按现有题面,极限为
    \[
        \boxed{\ln b\,\sin1\,\sin\sqrt b}.
    \]
    由于求和项完全不依赖 $i$,原题很可能漏写了与 $i$ 有关的指数或因子,需要对照原始资料核对.
\end{solution}

%例题6.1.11
\begin{example}
    求极限
    \[
        \lim\limits_{x\to+\infty}\dfrac{1}{x}\int_0^x|\sin t|\,\mathrm{d}t.
    \]
\end{example}
\exampleblank

\begin{solution}
    写成 $x=k\pi+r$,其中 $k\in\mathbb{N}$,$0\leq r<\pi$.由于
    \[
        \int_0^\pi|\sin t|\,\mathrm{d}t=2,
    \]
    故
    \[
        \int_0^x|\sin t|\,\mathrm{d}t=2k+O(1).
    \]
    又 $x=k\pi+O(1)$,所以
    \[
        \lim\limits_{x\to+\infty}\dfrac{1}{x}
        \int_0^x|\sin t|\,\mathrm{d}t
        =\boxed{\dfrac{2}{\pi}}.
    \]
\end{solution}

%例题6.1.12
\begin{example}
    求证:
    \begin{enumerate}
        \item $\displaystyle\lim\limits_{n\to\infty}\int_0^{\frac{\pi}{2}}\sin^n x\,\mathrm{d}x=0$;
        \item $\displaystyle\lim\limits_{n\to\infty}\int_0^{\frac{\pi}{2}}(1-\sin x)^n\,\mathrm{d}x=0$,其中 $a\in(0,1)$;
        \item $\displaystyle\lim\limits_{n\to\infty}\int_0^1e^{x^n}\,\mathrm{d}x=1$.
    \end{enumerate}
\end{example}
\exampleblank

\begin{solution}
    \begin{enumerate}
        \item 在 $0\leq x<\dfrac{\pi}{2}$ 上有 $\sin^n x\to0$,且 $0\leq\sin^n x\leq1$.由控制收敛定理,积分趋于 $0$.

        \item 在 $0<x\leq\dfrac{\pi}{2}$ 上有 $(1-\sin x)^n\to0$,且该函数被 $1$ 控制,故积分也趋于 $0$.题面末尾的``其中 $a\in(0,1)$''未在公式中出现,应为多余文字或漏项.

        \item 对 $0\leq x<1$,有 $x^n\to0$,故 $e^{x^n}\to1$;同时 $1\leq e^{x^n}\leq e$.由控制收敛定理,
              \[
                  \lim\limits_{n\to\infty}\int_0^1e^{x^n}\,\mathrm{d}x
                  =\int_0^1 1\,\mathrm{d}x=1.
              \]
    \end{enumerate}
\end{solution}

%例题6.1.13
\begin{example}
    设 $f(x)$ 严格单调递减,在 $\lbrack 0,1\rbrack$ 上连续,$f(0)=1,f(1)=0$.试证明:对任意 $\delta\in(0,1)$ 有
    \begin{enumerate}
        \item
              \[
                  \lim\limits_{n\to\infty}
                  \dfrac{\displaystyle\int_\delta^1[f(x)]^n\,\mathrm{d}x}
                  {\displaystyle\int_0^\delta[f(x)]^n\,\mathrm{d}x}=0;
              \]
        \item
              \[
                  \lim\limits_{n\to\infty}
                  \dfrac{\displaystyle\int_0^\delta[f(x)]^{n+1}\,\mathrm{d}x}
                  {\displaystyle\int_0^1[f(x)]^n\,\mathrm{d}x}=1.
              \]
    \end{enumerate}
\end{example}
\exampleblank

\begin{solution}
    \begin{enumerate}
        \item 固定 $c\in(0,\delta)$.由严格递减性,
              \[
                  \int_\delta^1[f(x)]^n\,\mathrm{d}x
                  \leq(1-\delta)[f(\delta)]^n,
              \]
              而
              \[
                  \int_0^\delta[f(x)]^n\,\mathrm{d}x
                  \geq c[f(c)]^n.
              \]
              因 $f(\delta)<f(c)$,两式之比趋于零.

        \item 记
              \[
                  D_n=\int_0^1[f(x)]^n\,\mathrm{d}x.
              \]
              则
              \[
                  \begin{aligned}
                      0\leq D_n-\int_0^\delta[f(x)]^{n+1}\,\mathrm{d}x
                      ={} & \int_0^\delta[f(x)]^n[1-f(x)]\,\mathrm{d}x \\
                          & +\int_\delta^1[f(x)]^n\,\mathrm{d}x.
                  \end{aligned}
              \]
              对任意 $\varepsilon\in(0,\delta)$,第一项在 $\lbrack0,\varepsilon\rbrack$ 上不超过
              \[
                  [1-f(\varepsilon)]D_n,
              \]
              而在 $\lbrack\varepsilon,\delta\rbrack$ 上以及第二项都由第 (1) 问的估计可知相对于 $D_n$ 趋于零.再令 $\varepsilon\to0^+$,利用 $f(\varepsilon)\to1$,得到
              \[
                  \dfrac{\displaystyle\int_0^\delta[f(x)]^{n+1}\,\mathrm{d}x}{D_n}
                  \longrightarrow1.
              \]
    \end{enumerate}
\end{solution}

%例题6.1.14
\begin{example}[\markstar]
    设 $f(x)\geq0,g(x)>0$,两函数在 $\lbrack a,b\rbrack$ 上连续,求证:
    \[
        \lim\limits_{n\to\infty}
        \left[\int_a^b[f(x)]^n g(x)\,\mathrm{d}x\right]^{\frac{1}{n}}
        =\max\limits_{a\leq x\leq b}f(x).
    \]
\end{example}
\exampleblank

\begin{solution}
    记
    \[
        M=\max\limits_{a\leq x\leq b}f(x),
        \qquad
        G=\int_a^b g(x)\,\mathrm{d}x>0.
    \]
    上界为
    \[
        \left[\int_a^b f(x)^ng(x)\,\mathrm{d}x\right]^{\frac{1}{n}}
        \leq M G^{\frac{1}{n}}\longrightarrow M.
    \]
    任取 $\varepsilon>0$.由 $f$ 的连续性,在某个正长度闭区间 $J\subseteq\lbrack a,b\rbrack$ 上有 $f(x)\geq M-\varepsilon$.又因 $g>0$ 且连续,
    \[
        \int_Jg(x)\,\mathrm{d}x>0.
    \]
    于是
    \[
        \left[\int_a^b f(x)^ng(x)\,\mathrm{d}x\right]^{\frac{1}{n}}
        \geq(M-\varepsilon)
        \left[\int_Jg(x)\,\mathrm{d}x\right]^{\frac{1}{n}}
        \longrightarrow M-\varepsilon.
    \]
    令 $\varepsilon\to0^+$,由夹逼定理得到结论.
\end{solution}

拟合法:当证明极限 $\displaystyle\lim\limits_{n\to\infty}a_n=a$ 时,若将 $a$ 写成类似 $a_n$ 的形式,则通过证明 $\displaystyle\lim\limits_{n\to\infty}(a_n-a)=0$,将会大大简化问题.

%例题6.1.15
\begin{example}
    设 $f(x)$ 在 $\lbrack 0,1\rbrack$ 上连续,试证:
    \[
        \lim\limits_{h\to0^+}\int_0^1\dfrac{h}{h^2+x^2}f(x)\,\mathrm{d}x
        =\dfrac{\pi}{2}f(0).
    \]
\end{example}
\exampleblank

\begin{solution}
    作代换 $x=ht$,原积分化为
    \[
        I_h=\int_0^{\frac{1}{h}}\dfrac{f(ht)}{1+t^2}\,\mathrm{d}t.
    \]
    写成
    \[
        I_h-f(0)\dfrac{\pi}{2}
        =\int_0^{\frac{1}{h}}\dfrac{f(ht)-f(0)}{1+t^2}\,\mathrm{d}t
        -f(0)\int_{\frac{1}{h}}^{+\infty}\dfrac{\mathrm{d}t}{1+t^2}.
    \]
    第二项显然趋于零.对第一项,先取 $R$ 足够大,使核函数在 $\lbrack R,+\infty)$ 上的积分任意小;再在固定区间 $\lbrack0,R\rbrack$ 上利用 $f(ht)\to f(0)$ 的一致性.由此第一项也趋于零,故
    \[
        I_h\longrightarrow\dfrac{\pi}{2}f(0).
    \]
\end{solution}

%例题6.1.16
\begin{example}
    设 $f$ 为连续函数,证明:
    \[
        \lim\limits_{n\to\infty}\dfrac{2}{\pi}
        \int_0^1\dfrac{n}{n^2x^2+1}f(x)\,\mathrm{d}x=f(0).
    \]
\end{example}
\exampleblank

\begin{solution}
    令 $t=nx$,则左端积分为
    \[
        \dfrac{2}{\pi}\int_0^n
        \dfrac{f\left(\frac{t}{n}\right)}{1+t^2}\,\mathrm{d}t.
    \]
    与上一题完全相同,先控制核函数在大 $t$ 处的尾积分,再在固定有界区间上使用 $f\left(\dfrac{t}{n}\right)\to f(0)$ 的一致性,得到
    \[
        \lim\limits_{n\to\infty}\dfrac{2}{\pi}
        \int_0^n\dfrac{f\left(\frac{t}{n}\right)}{1+t^2}\,\mathrm{d}t
        =\dfrac{2}{\pi}f(0)\int_0^{+\infty}
        \dfrac{\mathrm{d}t}{1+t^2}
        =f(0).
    \]
\end{solution}

%例题6.1.17
\begin{example}
    设 $f(x)$ 在 $\lbrack a,b\rbrack$ 上可积,在 $x=b$ 处左连续,求证:
    \[
        \lim\limits_{n\to\infty}
        \dfrac{n+1}{(b-a)^{n+1}}
        \int_a^b(x-a)^n f(x)\,\mathrm{d}x=f(b).
    \]
\end{example}
\exampleblank

\begin{solution}
    令 $x=a+(b-a)t$,原式化为
    \[
        (n+1)\int_0^1t^n f(a+(b-a)t)\,\mathrm{d}t.
    \]
    核函数 $(n+1)t^n$ 非负且在 $\lbrack0,1\rbrack$ 上积分为 $1$.任取 $\delta\in(0,1)$,在 $\lbrack0,1-\delta\rbrack$ 上有
    \[
        (n+1)\int_0^{1-\delta}t^n\,\mathrm{d}t
        =(1-\delta)^{n+1}\longrightarrow0.
    \]
    而在 $t$ 接近 $1$ 时,由 $f$ 在 $b$ 处左连续,
    \[
        f(a+(b-a)t)\longrightarrow f(b).
    \]
    将积分分成上述两段并用夹逼估计,即得极限为 $f(b)$.
\end{solution}

%例题6.1.18
\begin{example}
    设 $f\in C\lbrack A,B\rbrack$,$A<a<b<B$,试证:
    \[
        \lim\limits_{h\to0}\int_a^b\dfrac{f(x+h)-f(x)}{h}\,\mathrm{d}x
        =f(b)-f(a).
    \]
\end{example}
\exampleblank

\begin{solution}
    对充分小的 $h$,由变量平移,
    \[
        \begin{aligned}
            \int_a^b\dfrac{f(x+h)-f(x)}{h}\,\mathrm{d}x
             & =
            \dfrac{1}{h}\left(
            \int_b^{b+h}f(u)\,\mathrm{d}u
            -\int_a^{a+h}f(u)\,\mathrm{d}u\right).
        \end{aligned}
    \]
    当 $h\to0$ 时,两个积分平均值分别趋于 $f(b)$ 与 $f(a)$,故原极限为 $f(b)-f(a)$.
\end{solution}

%例题6.1.19
\begin{example}
    设 $f(x)$ 在 $\lbrack a,b\rbrack$ 上可导,$f''(x)$ 在 $\lbrack a,b\rbrack$ 上可积.对任意 $n\in\mathbb{N}$,记
    \[
        A_n=\sum\limits_{i=1}^n f\left(a+i\dfrac{b-a}{n}\right)\dfrac{b-a}{n}
        -\int_a^b f(x)\,\mathrm{d}x.
    \]
    试证:
    \[
        \lim\limits_{n\to\infty}nA_n
        =\dfrac{b-a}{2}[f(b)-f(a)].
    \]
\end{example}
\exampleblank

\begin{solution}
    记
    \[
        h=\dfrac{b-a}{n},
        \qquad x_i=a+ih.
    \]
    逐段写出右端矩形公式的误差:
    \[
        \begin{aligned}
            A_n
             & =\sum\limits_{i=1}^{n}
            \int_{x_{i-1}}^{x_i}[f(x_i)-f(x)]\,\mathrm{d}x \\
             & =\sum\limits_{i=1}^{n}
            \int_{x_{i-1}}^{x_i}(t-x_{i-1})f'(t)\,\mathrm{d}t.
        \end{aligned}
    \]
    由于 $f'$ 一致连续,
    \[
        A_n=\dfrac{h^2}{2}
        \sum\limits_{i=1}^{n}f'(x_i)+o(h).
    \]
    于是
    \[
        \begin{aligned}
            nA_n
             & =\dfrac{b-a}{2}
            \left[h\sum\limits_{i=1}^{n}f'(x_i)\right]+o(1) \\
             & \longrightarrow\dfrac{b-a}{2}
            \int_a^b f'(x)\,\mathrm{d}x                     \\
             & =\dfrac{b-a}{2}[f(b)-f(a)].
        \end{aligned}
    \]
\end{solution}

%例题6.1.20
\begin{example}[\markstar]
    设 $f(x)$ 在 $\lbrack a,b\rbrack$ 上二次可微,且 $f''(x)$ 在 $\lbrack a,b\rbrack$ 上可积,记
    \[
        B_n=\int_a^b f(x)\,\mathrm{d}x
        -\dfrac{b-a}{n}\sum\limits_{i=1}^n
        f\left[a+(2i-1)\dfrac{b-a}{2n}\right].
    \]
    试证:
    \[
        \lim\limits_{n\to\infty}n^2B_n
        =\dfrac{(b-a)^2}{24}[f'(b)-f'(a)].
    \]
\end{example}
\exampleblank

\begin{solution}
    记 $h=\dfrac{b-a}{n}$,并令 $m_i=a+\left(i-\dfrac{1}{2}\right)h$.在每个小区间上使用带积分余项的二阶 Taylor 公式,得到
    \[
        \int_{m_i-\frac{h}{2}}^{m_i+\frac{h}{2}}f(x)\,\mathrm{d}x
        -hf(m_i)
        =\dfrac{h^3}{24}f''(m_i)+o(h^3),
    \]
    其中各小区间余项求和后为 $o(h^2)$.因此
    \[
        B_n=\dfrac{h^2}{24}
        \left[h\sum\limits_{i=1}^{n}f''(m_i)\right]+o(h^2).
    \]
    乘以 $n^2$,并利用 $n^2h^2=(b-a)^2$ 以及中点 Riemann 和收敛,得到
    \[
        \begin{aligned}
            n^2B_n
             & \longrightarrow\dfrac{(b-a)^2}{24}
            \int_a^b f''(x)\,\mathrm{d}x          \\
             & =\dfrac{(b-a)^2}{24}[f'(b)-f'(a)].
        \end{aligned}
    \]
\end{solution}

\begin{remark}
    因此,若设
    \[
        A_n=\sum\limits_{k=1}^n\dfrac{1}{n+k},
        \qquad
        B_n=\sum\limits_{k=1}^n\dfrac{2}{2n+2k-1},
    \]
    则
    \[
        \lim\limits_{n\to\infty}n(\ln2-A_n)=\dfrac{1}{4},
        \qquad
        \lim\limits_{n\to\infty}n^2(\ln2-B_n)=\dfrac{1}{32}.
    \]
\end{remark}

%例题6.1.21
\begin{example}[\markstar]
    设 $f(x)$ 在 $\lbrack a,b\rbrack$ 上可积,$g(x)\geq0$,且 $g$ 是以 $T>0$ 为周期,在 $\lbrack 0,T\rbrack$ 上可积的函数.试证:
    \[
        \lim\limits_{n\to\infty}\int_a^b f(x)g(nx)\,\mathrm{d}x
        =\dfrac{1}{T}\int_0^Tg(x)\,\mathrm{d}x\int_a^b f(x)\,\mathrm{d}x.
    \]
\end{example}
\exampleblank

\begin{solution}
    记周期平均值
    \[
        \overline g=\dfrac{1}{T}\int_0^Tg(t)\,\mathrm{d}t.
    \]
    先对区间示性函数证明.任取 $u<v$,作代换 $t=nx$,并把 $\lbrack nu,nv\rbrack$ 分成完整周期与至多两个不足一周期的余段,可得
    \[
        \int_u^v g(nx)\,\mathrm{d}x
        =(v-u)\overline g+O\left(\dfrac{1}{n}\right).
    \]
    因此结论对阶梯函数 $f$ 成立.

    一般的 Riemann 可积函数可在 $L^1$ 意义下由阶梯函数逼近.又因周期可积函数 $g$ 有界,若 $s$ 是阶梯函数,则
    \[
        \left|\int_a^b[f(x)-s(x)]g(nx)\,\mathrm{d}x\right|
        \leq\|g\|_\infty\int_a^b|f(x)-s(x)|\,\mathrm{d}x.
    \]
    先选取 $s$ 使右端任意小,再令 $n\to\infty$,便得到
    \[
        \int_a^b f(x)g(nx)\,\mathrm{d}x
        \longrightarrow\overline g\int_a^bf(x)\,\mathrm{d}x.
    \]
\end{solution}

%例题6.1.22
\begin{example}[\markstar\ Riemann 引理]
    若 $f(x)$ 在 $\lbrack a,b\rbrack$ 上可积,$g(x)$ 以 $T$ 为周期,在 $\lbrack 0,T\rbrack$ 上可积,则
    \[
        \lim\limits_{n\to\infty}\int_a^b f(x)g(nx)\,\mathrm{d}x
        =\dfrac{1}{T}\int_0^Tg(x)\,\mathrm{d}x\int_a^b f(x)\,\mathrm{d}x.
    \]
\end{example}
\exampleblank

\begin{solution}
    证明与上一题相同,非负性并非必需.具体地,令
    \[
        \overline g=\dfrac{1}{T}\int_0^Tg(t)\,\mathrm{d}t.
    \]
    对任意区间 $\lbrack u,v\rbrack$,周期分割给出
    \[
        \int_u^v[g(nx)-\overline g]\,\mathrm{d}x
        =O\left(\dfrac{1}{n}\right).
    \]
    所以结论先对阶梯函数成立.再用阶梯函数在 $L^1$ 中逼近 Riemann 可积的 $f$,并利用 $g$ 的有界性控制逼近误差,得到
    \[
        \int_a^b f(x)[g(nx)-\overline g]\,\mathrm{d}x\longrightarrow0.
    \]
    展开 $\overline g$ 即为题设结论.
\end{solution}

%例题6.1.23
\begin{example}
    设函数 $f(x)$ 在区间 $\lbrack 0,1\rbrack$ 上有一阶连续导数,且 $f(0)=f(1)$,$g(x)$ 是周期为 $1$ 的连续函数,记
    \[
        a_n=\int_0^1 f(x)g(x)\,\mathrm{d}x,
    \]
    求证:
    \[
        \lim\limits_{n\to\infty}na_n=0.
    \]
\end{example}
\exampleblank

\begin{solution}
    本题当前定义
    \[
        a_n=\int_0^1f(x)g(x)\,\mathrm{d}x
    \]
    完全不含 $n$,所以结论一般不成立.取 $f(x)\equiv1$,$g(x)\equiv1$,则所有条件均满足,但 $a_n=1$,从而 $na_n=n\not\to0$.

    原题很可能应为
    \[
        a_n=\int_0^1f(x)g(nx)\,\mathrm{d}x
    \]
    并且还需令 $g$ 的周期平均值为零;即便如此,要得到 $na_n\to0$ 仍通常需要更强条件或对原式作分部积分.现有题面必须对照原始资料核对.
\end{solution}

%例题6.1.24
\begin{example}
    设 $s(x)=4[x]-2[2x]+1$,$f$ 在 $\lbrack 0,1\rbrack$ 上可积,证明:
    \[
        \lim\limits_{n\to\infty}\int_0^1 f(x)s(nx)\,\mathrm{d}x=0.
    \]
\end{example}
\exampleblank

\begin{solution}
    函数 $s$ 以 $1$ 为周期.在一个周期内,
    \[
        s(x)=
        \begin{cases}
            1,  & 0\leq x<\dfrac{1}{2}, \\
            -1, & \dfrac{1}{2}\leq x<1,
        \end{cases}
    \]
    故
    \[
        \int_0^1s(x)\,\mathrm{d}x=0.
    \]
    对上一题的 Riemann 引理取 $g=s$,$T=1$,便有
    \[
        \lim\limits_{n\to\infty}\int_0^1f(x)s(nx)\,\mathrm{d}x
        =\left(\int_0^1s(x)\,\mathrm{d}x\right)
        \left(\int_0^1f(x)\,\mathrm{d}x\right)=0.
    \]
\end{solution}


\newpage

\section{可积性}

\begin{theorem}[\markstar]
    $\lbrack a,b\rbrack$ 上的有界函数可积的充分必要条件是:对任意 $\varepsilon>0,\sigma>0$,存在划分
    \[
        \Delta:a=x_0<x_1<\cdots<x_n=b,
    \]
    其相应于 $\omega\geq\varepsilon$ 的子区间 $\Delta x_i$ 的长度总和小于 $\sigma$.
\end{theorem}

\begin{corollary}[\markstar]
    若定义在 $\lbrack a,b\rbrack$ 上的有界函数只有有限个不连续点,则 $f(x)$ 在 $\lbrack a,b\rbrack$ 上可积.
\end{corollary}

%例题6.2.1
\begin{example}
    试证明下列函数 $f(x)$ 在区间 $\lbrack 0,1\rbrack$ 上可积:
    \begin{enumerate}
        \item
              \[
                  f(x)=\begin{cases}
                      \sin\dfrac{1}{x}, & x\neq0, \\
                      0,                & x=0;
                  \end{cases}
              \]
        \item
              \[
                  f(x)=\begin{cases}
                      \dfrac{1}{x}-\dfrac{1}{\sin x}, & x\neq0, \\
                      0,                              & x=0;
                  \end{cases}
              \]
        \item
              \[
                  f(x)=\begin{cases}
                      \ln x\cdot\ln(1+x), & x\neq0, \\
                      0,                  & x=0.
                  \end{cases}
              \]
    \end{enumerate}
\end{example}
\exampleblank

\begin{solution}
    \begin{enumerate}
        \item 函数有界,并且在 $(0,1\rbrack$ 上连续;唯一可能的不连续点为 $0$.有界函数只有一个不连续点,故 Riemann 可积.

        \item 当 $x\to0$ 时,
              \[
                  \sin x=x-\dfrac{x^3}{6}+o(x^3),
                  \qquad
                  \dfrac{1}{\sin x}
                  =\dfrac{1}{x}+\dfrac{x}{6}+o(x).
              \]
              因而
              \[
                  \dfrac{1}{x}-\dfrac{1}{\sin x}\longrightarrow0.
              \]
              题设在 $0$ 处的取值使函数连续,故可积.

        \item 因 $\ln(1+x)\sim x$ 且 $x\ln x\to0$,有
              \[
                  \ln x\ln(1+x)\longrightarrow0
                  \qquad(x\to0^+).
              \]
              因此函数在 $0$ 处也连续,从而在闭区间上可积.
    \end{enumerate}
\end{solution}

\begin{proposition}[\markstar]
    设 $f(x)$ 为 $\lbrack a,b\rbrack$ 上的有界函数,$\{a_n\}\subseteq\lbrack a,b\rbrack$,且 $\displaystyle\lim\limits_{n\to\infty}a_n=c$.证明:若 $f(x)$ 在 $\lbrack a,b\rbrack$ 上只有 $a_n$($n=1,2,3,\ldots$)为其间断点,则 $f(x)$ 在 $\lbrack a,b\rbrack$ 上可积.
\end{proposition}

%例题6.2.2
\begin{example}
    试证明下列函数在 $\lbrack 0,1\rbrack$ 上是可积的:
    \begin{enumerate}
        \item
              \[
                  f(x)=\begin{cases}
                      0,                                      & x=0,    \\
                      \dfrac{1}{x}-\left[\dfrac{1}{x}\right], & x\neq0;
                  \end{cases}
              \]
        \item $f(x)=\operatorname{sgn}\left(\sin\dfrac{\pi}{x}\right)$.
    \end{enumerate}
\end{example}
\exampleblank

\begin{solution}
    \begin{enumerate}
        \item 对 $x\neq0$,该函数是 $\dfrac{1}{x}$ 的小数部分,故始终位于 $\lbrack0,1)$.它只可能在
              \[
                  x=\dfrac{1}{n},
                  \qquad n=1,2,\ldots,
              \]
              以及 $x=0$ 处不连续.不连续点集可数,故由 Riemann 可积判别准则,函数可积.

        \item 该函数有界,只在 $\sin\dfrac{\pi}{x}=0$ 的点
              \[
                  x=\dfrac{1}{n}
              \]
              以及聚点 $0$ 处可能不连续.不连续点集仍为可数集,故函数 Riemann 可积.
    \end{enumerate}
\end{solution}

%例题6.2.3
\begin{example}[\markstar]
    证明:Riemann 函数在 $\lbrack 0,1\rbrack$ 上可积.
\end{example}
\exampleblank

\begin{solution}
    记 Riemann 函数为
    \[
        R(x)=
        \begin{cases}
            \dfrac{1}{q}, & x=\dfrac{p}{q}\text{ 为既约分数}, \\
            0,            & x\text{ 为无理数或端点}.
        \end{cases}
    \]
    任取 $\varepsilon>0$,选 $N$ 使 $\dfrac{1}{N}<\varepsilon$.分母不超过 $N$ 的既约分数在 $\lbrack0,1\rbrack$ 中只有有限个,可用总长度任意小的有限个区间覆盖它们.在其余点上有 $0\leq R(x)<\varepsilon$.

    因此可作分割,使覆盖上述有限点的小区间对 Darboux 上和的贡献任意小,而其余区间上的振幅不超过 $\varepsilon$.故上,下 Darboux 和之差可以任意小,$R$ 可积;并且由于每个区间都含无理数,下和为零,所以
    \[
        \int_0^1R(x)\,\mathrm{d}x=0.
    \]
\end{solution}

\begin{theorem}[\markstar\ 四则运算]
    若 $f(x),g(x)$ 在 $\lbrack a,b\rbrack$ 上可积,则:
    \begin{enumerate}
        \item $kf(x),f(x)\pm g(x),f(x)g(x)$ 在 $\lbrack a,b\rbrack$ 上也可积,但 $\dfrac{f(x)}{g(x)}$($g(x)\neq0$)也不一定可积;
        \item 若 $|f(x)|\geq m>0$,则 $\dfrac{1}{f(x)}$ 在 $\lbrack a,b\rbrack$ 上也可积;
        \item 若 $f(x)$ 在 $\lbrack a,b\rbrack$ 上连续且恒不为零,则 $\dfrac{1}{f(x)}$ 在 $\lbrack a,b\rbrack$ 上也可积.
    \end{enumerate}
\end{theorem}

%例题6.2.4
\begin{example}[\markstar\ 复合]
    设 $f(x)$ 在 $\lbrack a,b\rbrack$ 上连续,$\varphi(t)$ 在 $\lbrack \alpha,\beta\rbrack$ 上可积,且
    \[
        a\leq\varphi(t)\leq b,
        \qquad t\in\lbrack \alpha,\beta\rbrack,
    \]
    则 $f(\varphi(t))$ 在 $\lbrack \alpha,\beta\rbrack$ 上仍然可积.
\end{example}
\exampleblank

\begin{solution}
    Riemann 可积函数 $\varphi$ 有界,且其不连续点集为零测集.若 $\varphi$ 在 $t_0$ 处连续,则由 $f$ 在紧区间 $\lbrack a,b\rbrack$ 上连续可知,复合函数 $f\circ\varphi$ 也在 $t_0$ 处连续.因此 $f\circ\varphi$ 的不连续点集包含在 $\varphi$ 的不连续点集中,仍为零测集.

    又因 $a\leq\varphi(t)\leq b$,$f(\varphi(t))$ 有界.由 Lebesgue 的 Riemann 可积判别准则,$f(\varphi(t))$ 在 $\lbrack\alpha,\beta\rbrack$ 上 Riemann 可积.
\end{solution}

\begin{remark}
    可积函数的复合不一定是可积函数,例如 $f(x)=\operatorname{sgn}x$ 与 $\varphi(t)=R(t)$ 为 Riemann 函数.
\end{remark}

%例题6.2.5
\begin{example}[\markstar]
    讨论 $f(x),|f(x)|,f^2(x)$ 在 $\lbrack a,b\rbrack$ 上可积性的关系.
\end{example}
\exampleblank

\begin{solution}
    三者的关系如下:
    \[
        f\in R\lbrack a,b\rbrack
        \quad\Longrightarrow\quad
        |f|\in R\lbrack a,b\rbrack
        \quad\Longleftrightarrow\quad
        f^2\in R\lbrack a,b\rbrack.
    \]
    第一个推出关系来自绝对值函数的连续性;$|f|$ 可积推出 $f^2=|f|^2$ 可积,而 $f^2$ 可积时
    \[
        |f|=\sqrt{f^2}
    \]
    也可积.

    反向一般不成立.取
    \[
        f(x)=
        \begin{cases}
            1,  & x\in\mathbb{Q},    \\
            -1, & x\notin\mathbb{Q}.
        \end{cases}
    \]
    则 $|f|\equiv1$,$f^2\equiv1$ 都可积,但 $f$ 在每一点都不连续,因而不可积.
\end{solution}

\begin{proposition}[\markstar]
    设 $f\in R\lbrack a,b\rbrack$,证明:$f$ 的连续点在 $\lbrack a,b\rbrack$ 中稠密.
\end{proposition}

\begin{corollary}[\markstar]
    设非负函数 $f\in R\lbrack a,b\rbrack$,则
    \[
        \int_a^b f(x)\,\mathrm{d}x=0
        \quad\Longleftrightarrow\quad
        f\text{ 在连续点上恒为零}.
    \]
\end{corollary}

%例题6.2.6
\begin{example}
    设 $f(x)$ 在 $\lbrack a,b\rbrack$ 上可积,试证:对于 $\lbrack a,b\rbrack$ 上任意可积函数 $g(x)$,恒有
    \[
        \int_a^b g(x)f(x)\,\mathrm{d}x=0,
    \]
    则函数 $f(x)$ 在连续点上恒为零.
\end{example}
\exampleblank

\begin{solution}
    在题设中取 $g=f$,得到
    \[
        \int_a^b f(x)^2\,\mathrm{d}x=0.
    \]
    设 $x_0$ 是 $f$ 的连续点.若 $f(x_0)\neq0$,由连续性,存在 $x_0$ 的一个小邻域及常数 $c>0$,使该邻域内 $f(x)^2\geq c$.这将导致
    \[
        \int_a^b f(x)^2\,\mathrm{d}x>0,
    \]
    矛盾.故每个连续点都满足 $f(x_0)=0$.
\end{solution}

%例题6.2.7
\begin{example}
    设在 $\lbrack -1,1\rbrack$ 上的连续函数 $f(x)$ 满足如下条件:对 $\lbrack -1,1\rbrack$ 上的任意连续偶函数 $g(x)$,
    \[
        \int_{-1}^{1}f(x)g(x)\,\mathrm{d}x=0.
    \]
    试证:$f(x)$ 是 $\lbrack -1,1\rbrack$ 上的奇函数.
\end{example}
\exampleblank

\begin{solution}
    令
    \[
        h(x)=f(x)+f(-x).
    \]
    则 $h$ 是连续偶函数,故可在题设中取 $g=h$.于是
    \[
        0=\int_{-1}^{1}f(x)h(x)\,\mathrm{d}x.
    \]
    作代换 $x\mapsto-x$,同一积分也等于
    \[
        \int_{-1}^{1}f(-x)h(x)\,\mathrm{d}x.
    \]
    两式相加得到
    \[
        0=\int_{-1}^{1}[f(x)+f(-x)]h(x)\,\mathrm{d}x
        =\int_{-1}^{1}h(x)^2\,\mathrm{d}x.
    \]
    由 $h$ 连续,必有 $h\equiv0$,即 $f(-x)=-f(x)$,所以 $f$ 为奇函数.
\end{solution}

%例题6.2.8
\begin{example}[特殊函数逼近]
    设 $f\in R(a,b)$,证明:对任意 $\varepsilon>0$,存在函数 $g(x)$,使得
    \[
        \int_a^b|f(x)-g(x)|\,\mathrm{d}x<\varepsilon,
    \]
    其中的 $g$ 是:
    \begin{enumerate}
        \item 阶梯函数;
        \item 折线函数;
        \item 连续函数;
        \item 连续可微函数.
    \end{enumerate}
\end{example}
\exampleblank

\begin{solution}
    任取 $\varepsilon>0$.
    \begin{enumerate}
        \item 由 $f$ 的 Riemann 可积性,可取一个分割及各小区间中的取样点,使相应阶梯函数 $g_1$ 满足
              \[
                  \int_a^b|f(x)-g_1(x)|\,\mathrm{d}x<\varepsilon.
              \]

        \item 阶梯函数只有有限个跳跃点.在每个跳跃点两侧取互不相交,总长度充分小的区间,并在这些区间内用线段连接左右常值,得到折线函数 $g_2$.因各函数有界,可使
              \[
                  \int_a^b|g_1-g_2|\,\mathrm{d}x
              \]
              任意小.

        \item 折线函数本身连续,故第 (2) 问同时给出了连续函数逼近.

        \item 在折线函数的有限个折点附近,用光滑圆弧或三次 Hermite 多项式替换折角,并保持这些小邻域总长度充分小,即得连续可微函数 $g_4$,且
              \[
                  \int_a^b|g_2-g_4|\,\mathrm{d}x
              \]
              任意小.
    \end{enumerate}
    在每一步把误差控制为 $\dfrac{\varepsilon}{4}$,再用三角不等式,即可使最终各种 $g$ 都满足题设误差界.
\end{solution}


\newpage

\section{积分估计与积分不等式}


\subsection{Newton-Leibniz 公式与上限变动积分}

%例题6.3.1
\begin{example}
    举例说明:可积不一定有原函数,原函数存在也不一定可积.
\end{example}
\exampleblank

\begin{solution}
    \begin{enumerate}
        \item 阶跃函数
              \[
                  f(x)=
                  \begin{cases}
                      0, & x<0,   \\
                      1, & x\geq0
                  \end{cases}
              \]
              在任意包含 $0$ 的闭区间上 Riemann 可积,但它不可能有原函数,因为导函数具有 Darboux 介值性,而该函数有跳跃.

        \item 令
              \[
                  F(x)=
                  \begin{cases}
                      x^2\sin\dfrac{1}{x^2}, & x\neq0, \\
                      0,                     & x=0.
                  \end{cases}
              \]
              则 $F$ 在 $\lbrack0,1\rbrack$ 上处处可导,故 $f=F'$ 有原函数 $F$.但对 $x\neq0$,
              \[
                  f(x)=2x\sin\dfrac{1}{x^2}
                  -\dfrac{2}{x}\cos\dfrac{1}{x^2},
              \]
              它在 $0$ 附近无界,因而不 Riemann 可积.
    \end{enumerate}
\end{solution}

\begin{theorem}[\markstar\ 原函数存在充分条件]
    设 $f\in R(\lbrack a,b\rbrack)$,且在点 $x_0\in\lbrack a,b\rbrack$ 处连续,则其上限变动积分 $\displaystyle\int_a^x f(t)\,\mathrm{d}t$ 在点 $x=x_0$ 处可导,且其导数是 $f(x_0)$.
\end{theorem}

\begin{theorem}[\markstar]
    若 $f\in R(\lbrack a,b\rbrack)$,则 $\displaystyle\int_a^x f(t)\,\mathrm{d}t$ 在 $\lbrack a,b\rbrack$ 上一致连续.
\end{theorem}

%例题6.3.2
\begin{example}
    设 $f\in C(\lbrack 0,1\rbrack)$ 且 $f(x)\geq0$,若有
    \[
        f^2(x)\leq1+2\int_0^x f(t)\,\mathrm{d}t,
    \]
    则 $f(x)\leq1+x$.
\end{example}
\exampleblank

\begin{solution}
    令
    \[
        F(x)=1+2\int_0^x f(t)\,\mathrm{d}t,
        \qquad
        y(x)=\sqrt{F(x)}.
    \]
    题设给出 $0\leq f(x)\leq y(x)$.又
    \[
        y'(x)=\dfrac{f(x)}{y(x)}\leq1,
        \qquad y(0)=1.
    \]
    因此 $y(x)\leq1+x$,进而
    \[
        f(x)\leq y(x)\leq1+x.
    \]
\end{solution}

%例题6.3.3
\begin{example}
    \begin{enumerate}
        \item 设 $f(x)$ 在 $\lbrack a,b\rbrack$ 上为正值可积函数,则存在 $\xi\in(a,b)$,使得
              \[
                  \int_a^\xi f(x)\,\mathrm{d}x
                  =\int_\xi^b f(x)\,\mathrm{d}x
                  =\dfrac{1}{2}\int_a^b f(x)\,\mathrm{d}x;
              \]
        \item 设 $f\in C(\lbrack a,b\rbrack)$ 且 $\displaystyle\int_a^b f(x)\,\mathrm{d}x=0$,则存在 $\xi\in(a,b)$,使得
              \[
                  \int_a^\xi f(x)\,\mathrm{d}x=f(\xi).
              \]
    \end{enumerate}
\end{example}
\exampleblank

\begin{solution}
    \begin{enumerate}
        \item 令
              \[
                  F(x)=\int_a^xf(t)\,\mathrm{d}t.
              \]
              因 $f>0$,$F$ 连续严格递增,且 $F(a)=0$,$F(b)=\displaystyle\int_a^bf$.由介值定理,存在唯一的 $\xi\in(a,b)$ 使
              \[
                  F(\xi)=\dfrac{1}{2}F(b).
              \]
              这正是题设两式.

        \item 令
              \[
                  G(x)=e^{-x}\int_a^xf(t)\,\mathrm{d}t.
              \]
              由积分为零,$G(a)=G(b)=0$.Rolle 定理给出某个 $\xi\in(a,b)$ 使 $G'(\xi)=0$.而
              \[
                  G'(x)=e^{-x}\left[f(x)-\int_a^xf(t)\,\mathrm{d}t\right],
              \]
              故 $\displaystyle\int_a^\xi f(t)\,\mathrm{d}t=f(\xi)$.
    \end{enumerate}
\end{solution}

%例题6.3.4
\begin{example}
    \begin{enumerate}
        \item 设 $f(x)$ 在 $\lbrack 0,1\rbrack$ 上可微,$0\leq f'(x)\leq1$ 且 $f(0)=0$,则
              \[
                  \left(\int_0^1f(x)\,\mathrm{d}x\right)^2
                  \geq\int_0^1f^3(x)\,\mathrm{d}x;
              \]
        \item 设 $f(x),g(x)$ 是 $\lbrack a,b\rbrack$ 上的递增连续函数,则
              \[
                  \int_a^bf(x)\,\mathrm{d}x\int_a^bg(x)\,\mathrm{d}x
                  \leq(b-a)\int_a^bf(x)g(x)\,\mathrm{d}x.
              \]
    \end{enumerate}
\end{example}
\exampleblank

\begin{solution}
    \begin{enumerate}
        \item 由 $f(0)=0$ 及 $0\leq f'\leq1$,有 $f\geq0$,并且
              \[
                  [f(x)^2]'=2f(x)f'(x)\leq2f(x).
              \]
              从 $0$ 到 $x$ 积分得
              \[
                  f(x)^2\leq2\int_0^xf(t)\,\mathrm{d}t.
              \]
              两端乘以 $f(x)\geq0$ 后再积分:
              \[
                  \begin{aligned}
                      \int_0^1f(x)^3\,\mathrm{d}x
                       & \leq2\int_0^1f(x)
                      \left[\int_0^xf(t)\,\mathrm{d}t\right]\mathrm{d}x \\
                       & =\left(\int_0^1f(x)\,\mathrm{d}x\right)^2.
                  \end{aligned}
              \]

        \item 直接计算差值:
              \[
                  \begin{aligned}
                       & (b-a)\int_a^bf(x)g(x)\,\mathrm{d}x
                      -\int_a^bf(x)\,\mathrm{d}x\int_a^bg(x)\,\mathrm{d}x \\
                       & \qquad=\dfrac{1}{2}\int_a^b\int_a^b
                      [f(x)-f(y)][g(x)-g(y)]\,\mathrm{d}x\,\mathrm{d}y.
                  \end{aligned}
              \]
              因 $f,g$ 同为递增函数,二重积分的被积函数非负,故结论成立.
    \end{enumerate}
\end{solution}

%例题6.3.5
\begin{example}
    设 $f\in C(\lbrack 0,1\rbrack)$,若有 $\displaystyle\int_0^x f(t)\,\mathrm{d}t\geq f(x)\geq0$,则 $f(x)\equiv0$.
\end{example}
\exampleblank

\begin{solution}
    令
    \[
        F(x)=\int_0^xf(t)\,\mathrm{d}t.
    \]
    则 $F(0)=0$,$F'(x)=f(x)\geq0$,且题设给出 $F'(x)\leq F(x)$.于是
    \[
        [e^{-x}F(x)]'\leq0.
    \]
    另一方面,$e^{-x}F(x)\geq0$ 且在 $x=0$ 时为零,所以它只能恒为零.故 $F\equiv0$,进而 $f=F'\equiv0$.
\end{solution}

%例题6.3.6
\begin{example}
    设 $f\in C^{(1)}(\lbrack 0,a\rbrack)$,$f(0)=0$,则
    \[
        \int_0^a|f(x)f'(x)|\,\mathrm{d}x
        \leq\dfrac{a}{2}\int_0^a[f'(x)]^2\,\mathrm{d}x.
    \]
    若条件改为 $f(a)=f(0)=0$ 或 $f\left(\dfrac{a}{2}\right)=0$,不等式右边系数该如何变化?
\end{example}
\exampleblank

\begin{solution}
    记 $q(x)=|f'(x)|$.由 $f(0)=0$,
    \[
        |f(x)|\leq\int_0^xq(t)\,\mathrm{d}t.
    \]
    因此
    \[
        \begin{aligned}
            \int_0^a|f(x)f'(x)|\,\mathrm{d}x
             & \leq\int_0^aq(x)\left[\int_0^xq(t)\,\mathrm{d}t\right]\mathrm{d}x \\
             & =\dfrac{1}{2}\left(\int_0^aq(x)\,\mathrm{d}x\right)^2             \\
             & \leq\dfrac{a}{2}\int_0^a[f'(x)]^2\,\mathrm{d}x.
        \end{aligned}
    \]

    若 $f(0)=f(a)=0$,则
    \[
        |f(x)|\leq\min\left\{\int_0^xq,\int_x^aq\right\}.
    \]
    以 $Q(x)=\int_0^xq$ 作变量可得
    \[
        \int_0^a|ff'|
        \leq\dfrac{1}{4}\left(\int_0^aq\right)^2
        \leq\dfrac{a}{4}\int_0^a(f')^2.
    \]
    若 $f\left(\dfrac{a}{2}\right)=0$,分别在 $\lbrack0,\dfrac{a}{2}\rbrack$,$\lbrack \dfrac{a}{2},a\rbrack$ 使用第一种估计,每段长度为 $\dfrac{a}{2}$,同样得到系数 $\boxed{\dfrac{a}{4}}$.
\end{solution}

%例题6.3.7
\begin{example}[Gronwall 引理]
    设 $f(x),g(x),\varphi(x)$ 均为 $\lbrack a,b\rbrack$ 上的连续函数,且 $g(x)$ 单调递增,$\varphi(x)\geq0$,且
    \[
        f(x)\leq g(x)+\int_a^x\varphi(t)f(t)\,\mathrm{d}t.
    \]
    证明:对任意 $x\in\lbrack a,b\rbrack$,有
    \[
        f(x)\leq g(x)e^{\int_a^x\varphi(s)\,\mathrm{d}s}.
    \]
\end{example}
\exampleblank

\begin{solution}
    令
    \[
        F(x)=\int_a^x\varphi(t)f(t)\,\mathrm{d}t,
        \qquad
        \Phi(x)=\int_a^x\varphi(t)\,\mathrm{d}t.
    \]
    则 $F(a)=0$,且由题设
    \[
        F'(x)=\varphi(x)f(x)
        \leq\varphi(x)[g(x)+F(x)].
    \]
    因此
    \[
        [e^{-\Phi(x)}F(x)]'
        \leq e^{-\Phi(x)}\varphi(x)g(x).
    \]
    从 $a$ 积分到 $x$,并利用 $g(t)\leq g(x)$,得到
    \[
        \begin{aligned}
            F(x)
             & \leq e^{\Phi(x)}g(x)
            \int_a^xe^{-\Phi(t)}\varphi(t)\,\mathrm{d}t \\
             & =g(x)[e^{\Phi(x)}-1].
        \end{aligned}
    \]
    最后由 $f(x)\leq g(x)+F(x)$ 得
    \[
        f(x)\leq g(x)e^{\Phi(x)}.
    \]
\end{solution}

%例题6.3.8
\begin{example}
    设 $f(x)\in C^1\lbrack 0,a\rbrack$,证明:
    \[
        |f(0)|\leq\dfrac{1}{a}\int_0^a|f(x)|\,\mathrm{d}x
        +\int_0^a|f'(x)|\,\mathrm{d}x.
    \]
\end{example}
\exampleblank

\begin{solution}
    对任意 $x\in\lbrack0,a\rbrack$,
    \[
        f(0)=f(x)-\int_0^xf'(t)\,\mathrm{d}t.
    \]
    故
    \[
        |f(0)|\leq|f(x)|+\int_0^a|f'(t)|\,\mathrm{d}t.
    \]
    对 $x$ 从 $0$ 到 $a$ 积分,再除以 $a$,便得
    \[
        |f(0)|\leq\dfrac{1}{a}\int_0^a|f(x)|\,\mathrm{d}x
        +\int_0^a|f'(x)|\,\mathrm{d}x.
    \]
\end{solution}

%例题6.3.9
\begin{example}
    设 $f(x)\in C^1\lbrack 0,a\rbrack$,证明:
    \[
        \int_0^1|f(x)|\,\mathrm{d}x
        \leq\max\left\{\int_0^1|f'(x)|\,\mathrm{d}x+
        \left|\int_0^1f(x)\,\mathrm{d}x\right|\right\}.
    \]
\end{example}
\exampleblank

\begin{solution}
    题面右端的 $\max\{\cdots\}$ 只有一个元素,等同于去掉 $\max$.由积分中值定理,存在 $c\in\lbrack0,1\rbrack$ 使
    \[
        f(c)=\int_0^1f(x)\,\mathrm{d}x.
    \]
    于是对任意 $x\in\lbrack0,1\rbrack$,
    \[
        |f(x)|
        \leq|f(c)|+\left|\int_c^xf'(t)\,\mathrm{d}t\right|
        \leq\left|\int_0^1f(x)\,\mathrm{d}x\right|
        +\int_0^1|f'(x)|\,\mathrm{d}x.
    \]
    再对 $x$ 积分即得题设不等式.
\end{solution}



\subsection{各种变形}

所谓各种变形即包括代数变形,换元(对称性),分部积分,中值定理,Taylor 公式等.另外,被积函数放大,缩小,区间放大和缩小也是获得估计的重要方法.



\methodsection{(A) 放缩被积函数}

%例题6.3.10
\begin{example}
    证明:
    \begin{enumerate}
        \item
              \[
                  \sqrt{2}<\int_{\sqrt{2}-1}^{\sqrt{2}}e^{-x^2}\,\mathrm{d}x<\sqrt{2};
              \]
        \item
              \[
                  0<\dfrac{\pi}{2}-\int_0^{\frac{\pi}{2}}\dfrac{\sin x}{x}\,\mathrm{d}x
                  <\dfrac{\pi^3}{144};
              \]
        \item
              \[
                  \dfrac{2}{9}\pi^2
                  \leq\int_{\frac{\pi}{6}}^{\frac{\pi}{2}}\dfrac{2x}{\sin x}\,\mathrm{d}x
                  \leq\dfrac{4}{9}\pi^2.
              \]
    \end{enumerate}
\end{example}
\exampleblank

\begin{solution}
    \begin{enumerate}
        \item 当前题面把严格下界与严格上界都写成了 $\sqrt2$,即要求同一个积分同时大于且小于 $\sqrt2$,逻辑上不可能成立.至少有一个端点常数存在排印错误,需要核对原始资料.

        \item 写成
              \[
                  \dfrac{\pi}{2}-\int_0^{\frac{\pi}{2}}
                  \dfrac{\sin x}{x}\,\mathrm{d}x
                  =\int_0^{\frac{\pi}{2}}\dfrac{x-\sin x}{x}\,\mathrm{d}x.
              \]
              对 $x>0$,有 $0<x-\sin x<\dfrac{x^3}{6}$,故
              \[
                  0<\dfrac{\pi}{2}-\int_0^{\frac{\pi}{2}}
                  \dfrac{\sin x}{x}\,\mathrm{d}x
                  <\dfrac{1}{6}\int_0^{\frac{\pi}{2}}x^2\,\mathrm{d}x
                  =\dfrac{\pi^3}{144}.
              \]

        \item 因 $\sin x\leq1$,
              \[
                  \int_{\frac{\pi}{6}}^{\frac{\pi}{2}}
                  \dfrac{2x}{\sin x}\,\mathrm{d}x
                  \geq\int_{\frac{\pi}{6}}^{\frac{\pi}{2}}2x\,\mathrm{d}x
                  =\dfrac{2\pi^2}{9}.
              \]
              又因弦估计 $\sin x\geq\dfrac{2x}{\pi}$,
              \[
                  \int_{\frac{\pi}{6}}^{\frac{\pi}{2}}
                  \dfrac{2x}{\sin x}\,\mathrm{d}x
                  \leq\int_{\frac{\pi}{6}}^{\frac{\pi}{2}}\pi\,\mathrm{d}x
                  =\dfrac{\pi^2}{3}<\dfrac{4\pi^2}{9}.
              \]
    \end{enumerate}
\end{solution}



\methodsection{(B) 换元法}

%例题6.3.11
\begin{example}
    证明:
    \[
        \int_0^{\sqrt{2\pi}}\sin x^2\,\mathrm{d}x>0.
    \]
\end{example}
\exampleblank

\begin{solution}
    令 $t=x^2$,则
    \[
        \int_0^{\sqrt{2\pi}}\sin x^2\,\mathrm{d}x
        =\dfrac{1}{2}\int_0^{2\pi}\dfrac{\sin t}{\sqrt t}\,\mathrm{d}t.
    \]
    把后一个积分分成 $\lbrack0,\pi\rbrack$ 与 $\lbrack\pi,2\pi\rbrack$,并在第二段令 $t=u+\pi$,得到
    \[
        \dfrac{1}{2}\int_0^\pi
        \sin u\left(\dfrac{1}{\sqrt u}
        -\dfrac{1}{\sqrt{u+\pi}}\right)\mathrm{d}u>0.
    \]
\end{solution}

%例题6.3.12
\begin{example}
    $f\in C(\lbrack a,b\rbrack)$ 且严格递增,则
    \[
        \int_a^bxf(x)\,\mathrm{d}x
        >\dfrac{a+b}{2}\int_a^bf(x)\,\mathrm{d}x.
    \]
\end{example}
\exampleblank

\begin{solution}
    使用严格 Chebyshev 积分不等式.函数 $x$ 与 $f(x)$ 在 $\lbrack a,b\rbrack$ 上同向严格递增,故
    \[
        \dfrac{1}{b-a}\int_a^bxf(x)\,\mathrm{d}x
        >
        \left(\dfrac{1}{b-a}\int_a^bx\,\mathrm{d}x\right)
        \left(\dfrac{1}{b-a}\int_a^bf(x)\,\mathrm{d}x\right).
    \]
    代入 $\displaystyle\int_a^b x\,\mathrm{d}x=\dfrac{(b-a)(a+b)}{2}$ 并整理,即得结论.
\end{solution}

%例题6.3.13
\begin{example}
    证明:
    \[
        \int_0^{\frac{\pi}{2}}\dfrac{\cos x}{1+x^2}\,\mathrm{d}x
        >\int_0^{\frac{\pi}{2}}\dfrac{\sin x}{1+x^2}\,\mathrm{d}x.
    \]
\end{example}
\exampleblank

\begin{solution}
    令
    \[
        g(x)=\cos x-\sin x,
        \qquad
        h(x)=\dfrac{1}{1+x^2}.
    \]
    二者都在 $\lbrack0,\dfrac{\pi}{2}\rbrack$ 上严格递减,且
    \[
        \int_0^{\frac{\pi}{2}}g(x)\,\mathrm{d}x=0.
    \]
    由严格 Chebyshev 积分不等式,
    \[
        \int_0^{\frac{\pi}{2}}g(x)h(x)\,\mathrm{d}x>0.
    \]
    展开该式就是题设两个积分之差为正.
\end{solution}



\methodsection{(C) 分部积分}

\begin{proposition}[\markstar]
    若 $f(x),g(x)$ 在 $\lbrack a,b\rbrack$ 上有 $2n$ 阶连续导数,且
    \[
        f^{(i)}(a)=f^{(i)}(b)=g^{(i)}(a)=g^{(i)}(b)=0,
        \qquad i=0,1,2,\ldots,n-1,
    \]
    则
    \[
        \int_a^bf^{(2n)}(x)g(x)\,\mathrm{d}x
        =\int_a^bf(x)g^{(2n)}(x)\,\mathrm{d}x.
    \]
\end{proposition}

%例题6.3.14
\begin{example}
    设 $f$ 在 $\lbrack a,b\rbrack$ 上二阶可微,$f(a)=f(b)=0$,$|f''(x)|\leq M$,证明:
    \[
        \left|\int_a^bf(x)\,\mathrm{d}x\right|
        \leq\dfrac{M}{2}(b-a)^3.
    \]
\end{example}
\exampleblank

\begin{solution}
    二次 Hermite 插值余项给出:对每个 $x\in(a,b)$,存在 $\xi_x\in(a,b)$ 使
    \[
        f(x)=\dfrac{f''(\xi_x)}{2}(x-a)(x-b).
    \]
    因此
    \[
        |f(x)|\leq\dfrac{M}{2}(x-a)(b-x).
    \]
    积分得更强的估计
    \[
        \left|\int_a^bf(x)\,\mathrm{d}x\right|
        \leq\int_a^b|f(x)|\,\mathrm{d}x
        \leq\dfrac{M}{12}(b-a)^3
        \leq\dfrac{M}{2}(b-a)^3.
    \]
\end{solution}

%例题6.3.15
\begin{example}
    设 $f(x)$ 在 $\lbrack a,b\rbrack$ 上有 $2n$ 阶连续导数,$|f^{(2n)}(x)|\leq M$,且
    \[
        f^{(i)}(a)=f^{(i)}(b)=0,
        \qquad i=0,1,2,\ldots,n-1.
    \]
    试证:
    \[
        \left|\int_a^bf(x)\,\mathrm{d}x\right|
        \leq\dfrac{(n!)^2M}{(2n)!(2n+1)!}(b-a)^{2n+1}.
    \]
\end{example}
\exampleblank

\begin{solution}
    把端点条件看作各具有重数 $n$ 的 Hermite 插值条件.其余项公式说明:对每个 $x\in(a,b)$,存在 $\xi_x\in(a,b)$ 使
    \[
        f(x)=\dfrac{f^{(2n)}(\xi_x)}{(2n)!}
        (x-a)^n(x-b)^n.
    \]
    故
    \[
        |f(x)|\leq\dfrac{M}{(2n)!}(x-a)^n(b-x)^n.
    \]
    积分并令 $x=a+(b-a)t$,得到
    \[
        \begin{aligned}
            \left|\int_a^bf(x)\,\mathrm{d}x\right|
             & \leq\dfrac{M(b-a)^{2n+1}}{(2n)!}
            \int_0^1t^n(1-t)^n\,\mathrm{d}t                \\
             & =\dfrac{(n!)^2M}{(2n)!(2n+1)!}(b-a)^{2n+1}.
        \end{aligned}
    \]
\end{solution}

%例题6.3.16
\begin{example}
    设 $f'(x)$ 在 $\lbrack 0,2\pi\rbrack$ 上连续,且 $f'(x)\geq0$,则对任意 $n\in\mathbb{Z}^+$,有
    \[
        \left|\int_0^{2\pi}f(x)\sin nx\,\mathrm{d}x\right|
        \leq\dfrac{2[f(2\pi)-f(0)]}{n}.
    \]
\end{example}
\exampleblank

\begin{solution}
    分部积分得到
    \[
        \begin{aligned}
            \int_0^{2\pi}f(x)\sin nx\,\mathrm{d}x
             & =
            \dfrac{f(0)-f(2\pi)}{n}
            +\dfrac{1}{n}\int_0^{2\pi}f'(x)\cos nx\,\mathrm{d}x.
        \end{aligned}
    \]
    因 $f'\geq0$,
    \[
        \int_0^{2\pi}f'(x)\,\mathrm{d}x=f(2\pi)-f(0).
    \]
    对前式取绝对值便得
    \[
        \left|\int_0^{2\pi}f(x)\sin nx\,\mathrm{d}x\right|
        \leq\dfrac{2[f(2\pi)-f(0)]}{n}.
    \]
\end{solution}



\methodsection{(D) 中值定理}

%例题6.3.17
\begin{example}
    设函数 $f(x)$ 在 $\lbrack 2,4\rbrack$ 上有连续的一阶导数,且 $f(2)=f(4)=0$,试证:
    \[
        \max\limits_{2\leq x\leq4}f'(x)
        \geq\int_2^4f(x)\,\mathrm{d}x.
    \]
\end{example}
\exampleblank

\begin{solution}
    本题按当前题面不成立.令 $t=x-2$,并取
    \[
        f(x)=t-\dfrac{t^4}{8},
        \qquad 0\leq t\leq2.
    \]
    则 $f(2)=f(4)=0$,并且
    \[
        f'(x)=1-\dfrac{t^3}{2},
        \qquad
        \max\limits_{2\leq x\leq4}f'(x)=1.
    \]
    但
    \[
        \int_2^4f(x)\,\mathrm{d}x
        =\int_0^2\left(t-\dfrac{t^4}{8}\right)\mathrm{d}t
        =\dfrac{6}{5}>1.
    \]
    因此结论缺少条件或常数有误,需要核对原始资料.
\end{solution}

%例题6.3.18
\begin{example}
    设 $f(x)$ 在 $\lbrack 0,1\rbrack$ 上有二阶连续导数,证明:
    \[
        \max\limits_{x\in\lbrack 0,1\rbrack}|f'(x)|
        \leq|f(1)-f(0)|+\int_0^1|f''(x)|\,\mathrm{d}x.
    \]
\end{example}
\exampleblank

\begin{solution}
    由 Lagrange 中值定理,存在 $c\in(0,1)$ 使
    \[
        f'(c)=f(1)-f(0).
    \]
    对任意 $x\in\lbrack0,1\rbrack$,
    \[
        |f'(x)|
        \leq|f'(c)|+\left|\int_c^xf''(t)\,\mathrm{d}t\right|
        \leq|f(1)-f(0)|+\int_0^1|f''(t)|\,\mathrm{d}t.
    \]
    对 $x$ 取最大值即得结论.
\end{solution}

%例题6.3.19
\begin{example}
    已知 $f(x)$ 在 $\lbrack 0,1\rbrack$ 上有二阶连续导数,且 $f(1)=f(0)=0$,证明:
    \[
        \int_0^1|f'(x)|\,\mathrm{d}x
        \leq\int_0^1|f''(x)|\,\mathrm{d}x+2\int_0^1|f(x)|\,\mathrm{d}x.
    \]
\end{example}
\exampleblank

\begin{solution}
    由 $f(0)=f(1)$,Rolle 定理给出 $c\in(0,1)$ 使 $f'(c)=0$.因此对任意 $x$,
    \[
        |f'(x)|\leq\int_0^1|f''(t)|\,\mathrm{d}t.
    \]
    在 $x\in\lbrack0,1\rbrack$ 上积分,得到更强的估计
    \[
        \int_0^1|f'(x)|\,\mathrm{d}x
        \leq\int_0^1|f''(x)|\,\mathrm{d}x.
    \]
    再在右端增加非负项 $2\displaystyle\int_0^1|f(x)|\,\mathrm{d}x$,即得题设结论.
\end{solution}

%例题6.3.20
\begin{example}
    已知 $f(x)$ 在 $\lbrack 0,1\rbrack$ 上有直到二阶的连续导数,证明:
    \[
        \int_0^1|f'(x)|\,\mathrm{d}x
        \leq\int_0^1|f''(x)|\,\mathrm{d}x+8\int_0^1|f(x)|\,\mathrm{d}x.
    \]
\end{example}
\exampleblank

\begin{solution}
    先证明存在 $c\in\lbrack0,1\rbrack$ 使
    \[
        |f'(c)|\leq8\int_0^1|f(x)|\,\mathrm{d}x.
    \]
    记
    \[
        A=\int_0^1|f(x)|\,\mathrm{d}x.
    \]
    若结论不成立,则由导函数的 Darboux 性质,$f'$ 保持定号,且处处有
    \[
        |f'(x)|>8A.
    \]
    不妨设 $f'(x)>8A$.若 $f$ 在 $z\in\lbrack0,1\rbrack$ 处为零,则由中值定理,
    \[
        |f(x)|>8A|x-z|,
    \]
    从而
    \[
        A>8A\int_0^1|x-z|\,\mathrm{d}x
        =4A\left[z^2+(1-z)^2\right]\geq2A,
    \]
    矛盾.若 $f$ 无零点,则 $f$ 保持定号;当 $f>0$ 时,
    $f(x)>f(0)+8Ax>8Ax$,故 $A>4A$;当 $f<0$ 时,从右端点反向比较同样得到 $A>4A$.情形 $f'(x)<-8A$ 完全类似.因此所需的 $c$ 必存在.

    于是对任意 $x$,
    \[
        |f'(x)|\leq|f'(c)|+\int_0^1|f''(t)|\,\mathrm{d}t.
    \]
    积分后得到
    \[
        \int_0^1|f'(x)|\,\mathrm{d}x
        \leq\int_0^1|f''(x)|\,\mathrm{d}x
        +8\int_0^1|f(x)|\,\mathrm{d}x.
    \]
\end{solution}

%例题6.3.21
\begin{example}
    $f(x)$ 在 $\lbrack 0,1\rbrack$ 上有连续二阶导数,$f(0)=f(1)=0$,$f(x)\neq0$(或 $f''(x)<0$),证明:
    \[
        \int_0^1\left|\dfrac{f''(x)}{f(x)}\right|\,\mathrm{d}x\geq4.
    \]
\end{example}
\exampleblank

\begin{solution}
    因 $f$ 在 $(0,1)$ 内不为零且连续,它保持定号.不妨设 $f>0$.令 $c$ 为 $f$ 的最大值点,则 $f'(c)=0$,并记
    \[
        q(x)=\left|\dfrac{f''(x)}{f(x)}\right|.
    \]
    对 $0<x<c$,
    \[
        |f'(x)|=\left|\int_x^cf''(t)\,\mathrm{d}t\right|
        \leq\max f\int_x^cq(t)\,\mathrm{d}t.
    \]
    再从 $0$ 积分到 $c$,并用 $f(0)=0$,得到
    \[
        1\leq\int_0^c t q(t)\,\mathrm{d}t
        \leq c\int_0^cq(t)\,\mathrm{d}t.
    \]
    同理,
    \[
        1\leq(1-c)\int_c^1q(t)\,\mathrm{d}t.
    \]
    故
    \[
        \begin{aligned}
            \int_0^1q(t)\,\mathrm{d}t
             & \geq\dfrac{1}{c}+\dfrac{1}{1-c} \\
             & \geq4.
        \end{aligned}
    \]
    若 $f''<0$,则由端点为零可知 $f>0$ 于内部,仍适用上述证明.
\end{solution}

%例题6.3.22
\begin{example}
    设 $f(x)\in C^2\lbrack a,b\rbrack$,且 $|f''(x)|\leq M$,证明:
    \begin{enumerate}
        \item 若 $f\left(\dfrac{a+b}{2}\right)=0$,则
              \[
                  \left|\int_a^bf(x)\,\mathrm{d}x\right|
                  \leq\dfrac{M}{24}(b-a)^3;
              \]
        \item 若 $f(a)=f(b)=0$,则
              \[
                  \left|\int_a^bf(x)\,\mathrm{d}x\right|
                  \leq\dfrac{M}{12}(b-a)^3.
              \]
    \end{enumerate}
\end{example}
\exampleblank

\begin{solution}
    \begin{enumerate}
        \item 记 $m=\dfrac{a+b}{2}$.Taylor 公式给出
              \[
                  f(x)=f'(m)(x-m)+R(x),
                  \qquad
                  |R(x)|\leq\dfrac{M}{2}(x-m)^2.
              \]
              对称区间上一次项积分为零,故
              \[
                  \left|\int_a^bf(x)\,\mathrm{d}x\right|
                  \leq\dfrac{M}{2}\int_a^b(x-m)^2\,\mathrm{d}x
                  =\dfrac{M}{24}(b-a)^3.
              \]

        \item 由二次 Hermite 插值余项,
              \[
                  |f(x)|\leq\dfrac{M}{2}(x-a)(b-x).
              \]
              因此
              \[
                  \left|\int_a^bf(x)\,\mathrm{d}x\right|
                  \leq\dfrac{M}{2}\int_a^b(x-a)(b-x)\,\mathrm{d}x
                  =\dfrac{M}{12}(b-a)^3.
              \]
    \end{enumerate}
\end{solution}

%例题6.3.23
\begin{example}
    设 $f(x)\in C^2\lbrack a,b\rbrack$,满足
    \[
        \left|\int_a^bf(x)\,\mathrm{d}x\right|
        <\int_a^b|f(x)|\,\mathrm{d}x,
    \]
    设
    \[
        \max\limits_{x\in\lbrack a,b\rbrack}|f'(x)|=M_1,
        \qquad
        \max\limits_{x\in\lbrack a,b\rbrack}|f''(x)|=M_2.
    \]
    求证:
    \[
        \left|\int_a^bf(x)\,\mathrm{d}x\right|
        \leq\dfrac{M_1}{2}(b-a)^2+\dfrac{M_2}{6}(b-a)^3.
    \]
\end{example}
\exampleblank

\begin{solution}
    严格不等式说明 $f$ 在区间上变号,故存在 $c\in(a,b)$ 使 $f(c)=0$.由 Taylor 公式及导数上界,
    \[
        |f(x)|
        \leq M_1|x-c|+\dfrac{M_2}{2}|x-c|^2.
    \]
    于是
    \[
        \begin{aligned}
            \left|\int_a^bf(x)\,\mathrm{d}x\right|
             & \leq M_1\int_a^b|x-c|\,\mathrm{d}x
            +\dfrac{M_2}{2}\int_a^b|x-c|^2\,\mathrm{d}x \\
             & \leq\dfrac{M_1}{2}(b-a)^2
            +\dfrac{M_2}{6}(b-a)^3.
        \end{aligned}
    \]
\end{solution}



\subsection{其它问题}

%例题6.3.24
\begin{example}
    设 $f(x)$ 在 $\lbrack 0,1\rbrack$ 上连续,
    \[
        \int_0^1f(x)\,\mathrm{d}x=0,
        \quad
        \int_0^1xf(x)\,\mathrm{d}x=0,
        \quad\ldots,\quad
        \int_0^1x^{n-1}f(x)\,\mathrm{d}x=0,
    \]
    且
    \[
        \int_0^1x^n f(x)\,\mathrm{d}x=1.
    \]
    求证:在 $\lbrack 0,1\rbrack$ 的某一部分上
    \[
        |f(x)|\geq2^n(n+1).
    \]
\end{example}
\exampleblank

\begin{solution}
    取多项式
    \[
        P(x)=(2x-1)^n=2^nx^n+Q_{n-1}(x),
    \]
    其中 $\deg Q_{n-1}\leq n-1$.由题设各阶矩条件,
    \[
        \int_0^1P(x)f(x)\,\mathrm{d}x=2^n.
    \]
    若记 $M=\max\limits_{0\leq x\leq1}|f(x)|$,则
    \[
        \begin{aligned}
            2^n
             & \leq M\int_0^1|2x-1|^n\,\mathrm{d}x \\
             & =\dfrac{M}{n+1}.
        \end{aligned}
    \]
    故
    \[
        M\geq2^n(n+1).
    \]
    由于最大值能够在闭区间上取得,所述部分非空.
\end{solution}

%例题6.3.25
\begin{example}
    设 $f(x)$ 在 $\lbrack 0,\dfrac{\pi}{2}\rbrack$ 上连续,
    \[
        \int_0^{\frac{\pi}{2}}f(x)\sin x\,\mathrm{d}x
        =\int_0^{\frac{\pi}{2}}f(x)\cos x\,\mathrm{d}x=0,
    \]
    证明:$f(x)$ 在 $\left(0,\dfrac{\pi}{2}\right)$ 内至少有两个零点.
\end{example}
\exampleblank

\begin{solution}
    反设 $f$ 在开区间内少于两个零点.若 $f$ 不变号,则因 $\sin x+\cos x>0$,
    \[
        \int_0^{\frac{\pi}{2}}f(x)(\sin x+\cos x)\,\mathrm{d}x\neq0,
    \]
    这与两个已知积分之和为零矛盾.

    若 $f$ 只在一点 $c\in(0,\dfrac{\pi}{2})$ 变号,则适当把 $f$ 替换为 $-f$ 后,$f(x)$ 与
    \[
        \sin(x-c)=\sin x\cos c-\cos x\sin c
    \]
    处处同号,且乘积不恒为零.因此
    \[
        \int_0^{\frac{\pi}{2}}f(x)\sin(x-c)\,\mathrm{d}x\neq0.
    \]
    但将 $\sin(x-c)$ 展开后,该积分由题设应为零,仍矛盾.故 $f$ 至少有两个内点零点.
\end{solution}

%例题6.3.26
\begin{example}[Hadamard 定理]
    设 $f(x)$ 是 $\lbrack a,b\rbrack$ 上连续凸函数,证明:对任意 $x_1,x_2\in\lbrack a,b\rbrack$ 且 $x_1<x_2$,有
    \[
        f\left(\dfrac{x_1+x_2}{2}\right)
        \leq\dfrac{1}{x_2-x_1}\int_{x_1}^{x_2}f(t)\,\mathrm{d}t
        \leq\dfrac{f(x_1)+f(x_2)}{2}.
    \]
\end{example}
\exampleblank

\begin{solution}
    对任意 $t\in\lbrack x_1,x_2\rbrack$,由凸性,
    \[
        f\left(\dfrac{x_1+x_2}{2}\right)
        \leq\dfrac{f(t)+f(x_1+x_2-t)}{2}.
    \]
    从 $x_1$ 到 $x_2$ 积分,得到左侧不等式.

    另一方面,凸函数图像不高于端点弦:
    \[
        f(t)\leq
        \dfrac{x_2-t}{x_2-x_1}f(x_1)
        +\dfrac{t-x_1}{x_2-x_1}f(x_2).
    \]
    再次积分并除以 $x_2-x_1$,得到右侧不等式.
\end{solution}

%例题6.3.27
\begin{example}
    设 $f(x)$ 是 $\lbrack 0,+\infty)$ 上的凸函数,求证:
    \[
        F(x)=\dfrac{1}{x}\int_0^xf(t)\,\mathrm{d}t
    \]
    为 $(0,+\infty)$ 上的凸函数.
\end{example}
\exampleblank

\begin{solution}
    作代换 $t=xs$,得到
    \[
        F(x)=\int_0^1f(xs)\,\mathrm{d}s.
    \]
    对每个固定的 $s\in\lbrack0,1\rbrack$,函数 $x\mapsto f(xs)$ 是凸函数.非负权积分保持凸性,故它们关于 $s$ 的积分 $F(x)$ 也是 $(0,+\infty)$ 上的凸函数.
\end{solution}

%例题6.3.28
\begin{example}[\markstar]
    设 $f(x),g(x)$ 在区间 $\lbrack a,b\rbrack$ 上连续,$p(x)\geq0$,$\displaystyle\int_a^bp(x)\,\mathrm{d}x>0$,且
    \[
        m\leq f(x)\leq M.
    \]
    $\varphi(x)$ 在 $\lbrack m,M\rbrack$ 上有定义,并有二阶导数,$\varphi''(x)>0$.求证:
    \[
        \varphi\left(
        \dfrac{\displaystyle\int_a^bp(x)f(x)\,\mathrm{d}x}
            {\displaystyle\int_a^bp(x)\,\mathrm{d}x}
        \right)
        \leq
        \dfrac{\displaystyle\int_a^bp(x)\varphi(f(x))\,\mathrm{d}x}
        {\displaystyle\int_a^bp(x)\,\mathrm{d}x}.
    \]
\end{example}
\exampleblank

\begin{solution}
    记
    \[
        P=\int_a^bp(x)\,\mathrm{d}x,
        \qquad
        \mu=\dfrac{1}{P}\int_a^bp(x)f(x)\,\mathrm{d}x.
    \]
    由 $m\leq f\leq M$,有 $\mu\in\lbrack m,M\rbrack$.凸函数在 $\mu$ 处的切线位于图像下方:
    \[
        \varphi(y)\geq\varphi(\mu)+\varphi'(\mu)(y-\mu).
    \]
    令 $y=f(x)$,乘以 $p(x)\geq0$ 后积分,得到
    \[
        \int_a^bp(x)\varphi(f(x))\,\mathrm{d}x
        \geq P\varphi(\mu),
    \]
    因为一次项的积分为零.两端除以 $P$ 即得 Jensen 不等式.
\end{solution}



\subsection{著名不等式}

\begin{theorem}[\marktwostar\ Schwarz 不等式]
    设 $f,g\in R\lbrack a,b\rbrack$,则
    \[
        \left(\int_a^bf(x)g(x)\,\mathrm{d}x\right)^2
        \leq\int_a^bf^2(x)\,\mathrm{d}x\int_a^bg^2(x)\,\mathrm{d}x.
    \]
\end{theorem}

\begin{theorem}[\markstar\ Young 不等式]
    设 $f$ 在 $\lbrack 0,+\infty)$ 上连续可微且严格单调递增,$f(0)=0$,$a,b>0$,则
    \[
        ab\leq\int_0^af(x)\,\mathrm{d}x+\int_0^bg(y)\,\mathrm{d}y,
    \]
    其中 $g(y)$ 是 $f(x)$ 的反函数,而等号当且仅当 $b=f(a)$ 时成立.
\end{theorem}

\begin{theorem}[\markstar\ Holder 不等式]
    设 $f,g\in R\lbrack a,b\rbrack$,$p,q$ 为满足 $\dfrac{1}{p}+\dfrac{1}{q}=1$ 的一对正实数,则
    \[
        \int_a^b|f(x)g(x)|\,\mathrm{d}x
        \leq
        \left(\int_a^b|f(x)|^p\,\mathrm{d}x\right)^{\frac{1}{p}}
        \left(\int_a^b|g(x)|^q\,\mathrm{d}x\right)^{\frac{1}{q}}.
    \]
\end{theorem}

\begin{theorem}[\markstar\ Minkowski 不等式]
    设 $f,g\in R\lbrack a,b\rbrack$,$1\leq p\leq+\infty$,则
    \[
        \left(\int_a^b|f(x)+g(x)|^p\,\mathrm{d}x\right)^{\frac{1}{p}}
        \leq
        \left(\int_a^b|f(x)|^p\,\mathrm{d}x\right)^{\frac{1}{p}}
        +\left(\int_a^b|g(x)|^p\,\mathrm{d}x\right)^{\frac{1}{p}}.
    \]
\end{theorem}

%例题6.3.29
\begin{example}
    设 $f\in C\lbrack 0,1\rbrack$,$0\leq f(x)<1$,则
    \[
        \int_0^1\dfrac{f(x)}{1-f(x)}\,\mathrm{d}x
        \geq
        \dfrac{\displaystyle\int_0^1\sqrt{f(x)}\,\mathrm{d}x}
        {1-\displaystyle\int_0^1f(x)\,\mathrm{d}x}.
    \]
\end{example}
\exampleblank

\begin{solution}
    本题按当前题面不成立.取常值函数 $f(x)\equiv\dfrac{1}{4}$,则
    \[
        \int_0^1\dfrac{f(x)}{1-f(x)}\,\mathrm{d}x
        =\dfrac{1}{3},
    \]
    而题设右端为
    \[
        \dfrac{\displaystyle\int_0^1\sqrt{f(x)}\,\mathrm{d}x}
        {1-\displaystyle\int_0^1f(x)\,\mathrm{d}x}
        =\dfrac{2}{3}.
    \]
    故不等式方向或右端分子(常见形式为该积分的平方)需要核对.
\end{solution}

%例题6.3.30
\begin{example}
    \begin{enumerate}
        \item 设 $f\in R\lbrack a,b\rbrack$,则
              \[
                  \left(\int_a^bf(x)\sin x\,\mathrm{d}x\right)^2
                  +\left(\int_a^bf(x)\cos x\,\mathrm{d}x\right)^2
                  \leq(b-a)\int_a^bf^2(x)\,\mathrm{d}x.
              \]
        \item 设 $f\in R\lbrack a,b\rbrack$ 且 $f(x)>0$,则
              \[
                  \left(\int_a^bf(x)\sin x\,\mathrm{d}x\right)^2
                  +\left(\int_a^bf(x)\cos x\,\mathrm{d}x\right)^2
                  \leq\left(\int_a^bf(x)\,\mathrm{d}x\right)^2.
              \]
    \end{enumerate}
\end{example}
\exampleblank

\begin{solution}
    \begin{enumerate}
        \item 记
              \[
                  A=\int_a^bf(x)\sin x\,\mathrm{d}x,
                  \qquad
                  B=\int_a^bf(x)\cos x\,\mathrm{d}x.
              \]
              选取 $\theta$ 使 $A\cos\theta+B\sin\theta=\sqrt{A^2+B^2}$.于是
              \[
                  \sqrt{A^2+B^2}
                  =\int_a^bf(x)\sin(x+\theta)\,\mathrm{d}x.
              \]
              由 Cauchy--Schwarz 不等式及 $\sin^2(x+\theta)\leq1$,
              \[
                  A^2+B^2\leq(b-a)\int_a^bf(x)^2\,\mathrm{d}x.
              \]

        \item 使用复数形式及三角不等式:
              \[
                  \begin{aligned}
                      \sqrt{A^2+B^2}
                       & =\left|\int_a^bf(x)e^{ix}\,\mathrm{d}x\right| \\
                       & \leq\int_a^bf(x)|e^{ix}|\,\mathrm{d}x
                      =\int_a^bf(x)\,\mathrm{d}x.
                  \end{aligned}
              \]
              两端平方即得结论.
    \end{enumerate}
\end{solution}

%例题6.3.31
\begin{example}
    设 $f\in R\lbrack a,b\rbrack$,若 $0<m\leq f(x)\leq M$,则:
    \begin{enumerate}
        \item
              \[
                  \left(\int_a^bf(x)\,\mathrm{d}x\right)
                  \left(\int_a^b\dfrac{\mathrm{d}x}{f(x)}\right)
                  \leq\dfrac{(m+M)^2}{4mM}(b-a)^2;
              \]
        \item
              \[
                  2\sqrt{\dfrac{m}{M}}(b-a)
                  \leq\dfrac{1}{M}\int_a^bf(x)\,\mathrm{d}x
                  +m\int_a^b\dfrac{\mathrm{d}x}{f(x)}.
              \]
    \end{enumerate}
\end{example}
\exampleblank

\begin{solution}
    \begin{enumerate}
        \item 由 $(M-f)(f-m)\geq0$,
              \[
                  f+\dfrac{mM}{f}\leq m+M.
              \]
              积分后记 $A=\displaystyle\int_a^bf$,$B=\displaystyle\int_a^b\dfrac{1}{f}$,得到
              \[
                  A+mMB\leq(m+M)(b-a).
              \]
              又由算术--几何平均不等式,
              \[
                  2\sqrt{mMAB}\leq A+mMB.
              \]
              合并并平方即得第 (1) 问.

        \item 逐点使用算术--几何平均不等式:
              \[
                  \dfrac{f(x)}{M}+\dfrac{m}{f(x)}
                  \geq2\sqrt{\dfrac{m}{M}}.
              \]
              在 $\lbrack a,b\rbrack$ 上积分即得结论.
    \end{enumerate}
\end{solution}

%例题6.3.32
\begin{example}
    设 $a>0$,证明:
    \[
        \int_0^\pi x^a\sin x\,\mathrm{d}x
        \int_0^{\frac{\pi}{2}}a^{-\cos x}\,\mathrm{d}x
        \geq\dfrac{\pi^3}{4}.
    \]
\end{example}
\exampleblank

\begin{solution}
    本题按当前题面不成立.取 $a=1$,则
    \[
        \int_0^\pi x^a\sin x\,\mathrm{d}x
        =\int_0^\pi x\sin x\,\mathrm{d}x=\pi,
    \]
    并且
    \[
        \int_0^{\frac{\pi}{2}}a^{-\cos x}\,\mathrm{d}x
        =\dfrac{\pi}{2}.
    \]
    两者乘积为 $\dfrac{\pi^2}{2}$,而
    \[
        \dfrac{\pi^2}{2}<\dfrac{\pi^3}{4}.
    \]
    因此右端常数或被积函数的底数,指数存在错误,需要核对原始资料.
\end{solution}

%例题6.3.33
\begin{example}
    设 $f(x)\in C^1\lbrack a,b\rbrack$,$f(a)=f(b)=0$,$\displaystyle\int_a^bf^2(x)\,\mathrm{d}x=1$,证明:
    \[
        \int_a^b[f'(x)]^2\,\mathrm{d}x
        \int_a^b[xf(x)]^2\,\mathrm{d}x
        \geq\dfrac{1}{4}.
    \]
\end{example}
\exampleblank

\begin{solution}
    对 $(xf(x)^2)'$ 积分.由 $f(a)=f(b)=0$,
    \[
        0=\int_a^b[xf(x)^2]'\,\mathrm{d}x
        =\int_a^bf(x)^2\,\mathrm{d}x
        +2\int_a^bxf(x)f'(x)\,\mathrm{d}x.
    \]
    结合 $\displaystyle\int_a^bf^2=1$,得到
    \[
        \left|\int_a^bxf(x)f'(x)\,\mathrm{d}x\right|
        =\dfrac{1}{2}.
    \]
    再用 Cauchy--Schwarz 不等式:
    \[
        \dfrac{1}{4}
        \leq\int_a^b[xf(x)]^2\,\mathrm{d}x
        \int_a^b[f'(x)]^2\,\mathrm{d}x.
    \]
\end{solution}

%例题6.3.34
\begin{example}
    设 $f(x)\in C^1\lbrack a,b\rbrack$,存在 $c\in\lbrack a,b\rbrack$,使得 $f(c)=0$.记
    \[
        M=\max\limits_{x\in\lbrack a,b\rbrack}|f(x)|,
    \]
    证明:
    \[
        M^2\leq(b-a)\int_a^b[f'(x)]^2\,\mathrm{d}x.
    \]
\end{example}
\exampleblank

\begin{solution}
    取 $x_0\in\lbrack a,b\rbrack$ 使 $|f(x_0)|=M$.由 $f(c)=0$,
    \[
        M=\left|\int_c^{x_0}f'(t)\,\mathrm{d}t\right|.
    \]
    使用 Cauchy--Schwarz 不等式,
    \[
        M^2
        \leq|x_0-c|\int_a^b[f'(t)]^2\,\mathrm{d}t
        \leq(b-a)\int_a^b[f'(t)]^2\,\mathrm{d}t.
    \]
\end{solution}

%例题6.3.35
\begin{example}
    设 $f(x)$ 在 $(0,1)$ 上可积,且 $\displaystyle\int_0^1xf(x)\,\mathrm{d}x=0$,证明:
    \[
        \int_0^1f^2(x)\,\mathrm{d}x
        \geq4\left(\int_0^1f(x)\,\mathrm{d}x\right)^2.
    \]
\end{example}
\exampleblank

\begin{solution}
    由 $\displaystyle\int_0^1xf(x)\,\mathrm{d}x=0$,对任意常数 $c$ 都有
    \[
        \int_0^1f(x)\,\mathrm{d}x
        =\int_0^1f(x)(1+cx)\,\mathrm{d}x.
    \]
    取 $c=-\dfrac{3}{2}$.由 Cauchy--Schwarz 不等式,
    \[
        \left(\int_0^1f(x)\,\mathrm{d}x\right)^2
        \leq\int_0^1f(x)^2\,\mathrm{d}x
        \int_0^1\left(1-\dfrac{3x}{2}\right)^2\mathrm{d}x.
    \]
    而最后一个积分等于 $\dfrac{1}{4}$,故
    \[
        \int_0^1f(x)^2\,\mathrm{d}x
        \geq4\left(\int_0^1f(x)\,\mathrm{d}x\right)^2.
    \]
\end{solution}


\newpage

\section{积分中值定理}

\begin{theorem}[\markstar\ 第一积分中值定理]
    设 $f(x)$ 是 $\lbrack a,b\rbrack$ 上的连续函数,则:
    \begin{enumerate}
        \item (强化形式)存在 $\xi\in(a,b)$,使得
              \[
                  \int_a^bf(x)\,\mathrm{d}x=f(\xi)(b-a).
              \]
        \item (推广形式)设 $g(x)$ 在 $\lbrack a,b\rbrack$ 上可积且不变号,则存在 $\xi\in(a,b)$,使得
              \[
                  \int_a^bf(x)g(x)\,\mathrm{d}x
                  =f(\xi)\int_a^bg(x)\,\mathrm{d}x.
              \]
    \end{enumerate}
\end{theorem}

\begin{remark}
    由证明过程可得,替代连续性,第一积分中值定理对具有介值性的可积函数仍然成立.
\end{remark}

%例题6.4.1
\begin{example}
    设 $f(x)$ 在 $(0,1)$ 上可微,且满足
    \[
        f(1)-2\int_0^1xf(x)\,\mathrm{d}x=0.
    \]
    求证:在 $(0,1)$ 内至少有一点 $\xi$,使得
    \[
        f'(\xi)=-\dfrac{f(\xi)}{3}.
    \]
\end{example}
\exampleblank

\begin{solution}
    本题按当前题面不成立.取非零常值函数 $f(x)\equiv C$,则
    \[
        f(1)-2\int_0^1xf(x)\,\mathrm{d}x
        =C-2C\int_0^1x\,\mathrm{d}x=0,
    \]
    所以题设条件成立;但 $f'(\xi)=0$,而 $-\dfrac{f(\xi)}{3}=-\dfrac{C}{3}\neq0$.因此原题还缺少条件,需要核对.
\end{solution}

%例题6.4.2
\begin{example}
    设 $f(x)$ 在 $\lbrack 0,1\rbrack$ 上连续,在 $(0,1)$ 内三次可微,并且满足条件:
    \begin{enumerate}
        \item $\displaystyle\int_0^{\frac{1}{4}}f(t)\,\mathrm{d}t=\dfrac{1}{4}f(0)$;
        \item $\displaystyle\int_{\frac{1}{4}}^1f(t)\,\mathrm{d}t=\dfrac{3}{4}f(1)$.
    \end{enumerate}
    试证:若 $f(0)=f(1)$,则存在 $\xi\in(0,1)$,使得 $f'''(\xi)=0$.
\end{example}
\exampleblank

\begin{solution}
    令 $C=f(0)=f(1)$.第一个积分条件等价于
    \[
        \int_0^{\frac{1}{4}}[f(t)-C]\,\mathrm{d}t=0,
    \]
    故由连续性,存在 $c_1\in(0,\dfrac{1}{4})$ 使 $f(c_1)=C$.同理,第二个条件给出某个 $c_2\in(\dfrac{1}{4},1)$,使 $f(c_2)=C$.

    于是函数 $f-C$ 在
    \[
        0,\quad c_1,\quad c_2,\quad1
    \]
    四个互不相同的点为零.连续使用三次 Rolle 定理,存在 $\xi\in(0,1)$ 使
    \[
        f'''(\xi)=0.
    \]
\end{solution}

%例题6.4.3
\begin{example}
    $\lbrack a,b\rbrack$ 上的连续函数序列 $\varphi_1,\varphi_2,\ldots,\varphi_n,\ldots$ 满足
    \[
        \int_a^b\varphi_n(x)\,\mathrm{d}x=1.
    \]
    证明:存在自然数 $N$ 及定数 $c_1,c_2,\ldots,c_N$,使得
    \[
        \sum\limits_{k=1}^N c_k^2=1,
        \qquad
        \max\limits_{x\in\lbrack a,b\rbrack}\left|\sum\limits_{k=1}^Nc_k\varphi_k(x)\right|>100.
    \]
\end{example}
\exampleblank

\begin{solution}
    取足够大的 $N$,使
    \[
        \dfrac{\sqrt N}{b-a}>100,
    \]
    并令
    \[
        c_1=\cdots=c_N=\dfrac{1}{\sqrt N}.
    \]
    显然 $\displaystyle\sum\limits_{k=1}^{N}c_k^2=1$.记
    \[
        F_N(x)=\sum\limits_{k=1}^{N}c_k\varphi_k(x).
    \]
    则
    \[
        \int_a^bF_N(x)\,\mathrm{d}x
        =\sum\limits_{k=1}^{N}c_k=\sqrt N.
    \]
    因此
    \[
        \max\limits_{x\in\lbrack a,b\rbrack}|F_N(x)|
        \geq\dfrac{1}{b-a}\left|\int_a^bF_N(x)\,\mathrm{d}x\right|
        =\dfrac{\sqrt N}{b-a}>100.
    \]
\end{solution}

%例题6.4.4
\begin{example}
    求
    \[
        \lim\limits_{x\to+\infty}
        \int_x^{x+2}t\left(\sin\dfrac{3}{t}\right)f(t)\,\mathrm{d}t,
    \]
    其中 $f(x)$ 可微,且 $\displaystyle\lim\limits_{t\to+\infty}f(t)=1$.
\end{example}
\exampleblank

\begin{solution}
    当 $t\to+\infty$ 时,
    \[
        t\sin\dfrac{3}{t}\longrightarrow3,
        \qquad
        f(t)\longrightarrow1.
    \]
    因此乘积在滑动区间 $\lbrack x,x+2\rbrack$ 上一致趋于 $3$.于是
    \[
        \int_x^{x+2}t\sin\dfrac{3}{t}f(t)\,\mathrm{d}t
        \longrightarrow\int_x^{x+2}3\,\mathrm{d}t
        =\boxed{6}.
    \]
\end{solution}

\begin{theorem}[\markstar\ 第二积分中值定理]
    \begin{enumerate}
        \item 在 $\lbrack a,b\rbrack$ 上,若 $f(x)$ 非负单调递减,$g(x)$ 可积,则存在 $\xi\in\lbrack a,b\rbrack$,使得
              \[
                  I=\int_a^bf(x)g(x)\,\mathrm{d}x
                  =f(a)\int_a^\xi g(x)\,\mathrm{d}x.
              \]
        \item 在 $\lbrack a,b\rbrack$ 上,若 $f(x)$ 非负单调递增,$g(x)$ 可积,则存在 $\xi\in\lbrack a,b\rbrack$,使得
              \[
                  I=\int_a^bf(x)g(x)\,\mathrm{d}x
                  =f(b)\int_\xi^bg(x)\,\mathrm{d}x.
              \]
        \item 在 $\lbrack a,b\rbrack$ 上,若 $f(x)$ 单调,$g(x)$ 可积,则存在 $\xi\in\lbrack a,b\rbrack$,使得
              \[
                  I=\int_a^bf(x)g(x)\,\mathrm{d}x
                  =f(a)\int_a^\xi g(x)\,\mathrm{d}x
                  +f(b)\int_\xi^bg(x)\,\mathrm{d}x.
              \]
    \end{enumerate}
\end{theorem}

\begin{remark}
    因此,若见积分问题条件中有单调性,可优先考虑第二积分中值定理.且由证明过程可知,第一式中的 $f(a)$ 可改写为 $f(a+0)$ 或比 $f(a+0)$ 大的任意常数 $A$,即存在 $\xi\in\lbrack a,b\rbrack$,使得
    \[
        I=\int_a^bf(x)g(x)\,\mathrm{d}x
        =A\int_a^\xi g(x)\,\mathrm{d}x.
    \]
    其中,$\xi$ 依赖于 $A$ 的选取.同理,第二式中的 $f(b)$ 可改写为 $f(b-0)$ 或比 $f(b-0)$ 大的任一常数 $B$.
\end{remark}

%例题6.4.5
\begin{example}
    证明:若 $f(x)$ 在 $\lbrack 0,1\rbrack$ 上严格单调递减,则:
    \begin{enumerate}
        \item 存在 $\theta\in(0,1)$,使得
              \[
                  \int_0^1f(x)\,\mathrm{d}x=\theta f(0)+(1-\theta)f(1);
              \]
        \item 对任意 $c>f(0)$,存在 $\theta\in(0,1)$,使得
              \[
                  \int_0^1f(x)\,\mathrm{d}x=\theta c+(1-\theta)f(1).
              \]
    \end{enumerate}
\end{example}
\exampleblank

\begin{solution}
    记
    \[
        I=\int_0^1f(x)\,\mathrm{d}x.
    \]
    因 $f$ 连续且严格递减,
    \[
        f(1)<I<f(0).
    \]
    \begin{enumerate}
        \item 取
              \[
                  \theta=\dfrac{I-f(1)}{f(0)-f(1)}.
              \]
              则 $0<\theta<1$,且 $I=\theta f(0)+(1-\theta)f(1)$.

        \item 对任意 $c>f(0)>I$,取
              \[
                  \theta=\dfrac{I-f(1)}{c-f(1)}.
              \]
              仍有 $0<\theta<1$,并且 $I=\theta c+(1-\theta)f(1)$.
    \end{enumerate}
\end{solution}

%例题6.4.6
\begin{example}
    设 $f:\lbrack a,b\rbrack\to\lbrack 0,1\rbrack$ 可导,$f'(x)$ 在 $\lbrack a,b\rbrack$ 上单调递减,$f'(b)>0$,求证:
    \[
        \int_a^b\cos f(x)\,\mathrm{d}x\leq\dfrac{2}{f'(b)}.
    \]
\end{example}
\exampleblank

\begin{solution}
    由 $f'$ 单调递减且 $f'(b)>0$,在全区间上都有
    \[
        f'(x)\geq f'(b)>0.
    \]
    故 $f$ 严格递增,可以作变量代换 $u=f(x)$.由于 $0\leq f(x)\leq1$,且 $\cos u>0$,有
    \[
        \begin{aligned}
            \int_a^b\cos f(x)\,\mathrm{d}x
             & =\int_{f(a)}^{f(b)}
            \dfrac{\cos u}{f'(f^{-1}(u))}\,\mathrm{d}u         \\
             & \leq\dfrac{1}{f'(b)}\int_0^1\cos u\,\mathrm{d}u \\
             & <\dfrac{1}{f'(b)}
            <\dfrac{2}{f'(b)}.
        \end{aligned}
    \]
\end{solution}

%例题6.4.7
\begin{example}
    \begin{enumerate}
        \item 设 $x>0,c>0$,证明:
              \[
                  \left|\int_x^{x-c}\sin t^2\,\mathrm{d}t\right|<\dfrac{1}{x};
              \]
        \item 设 $e^2<a<b$,证明:
              \[
                  \int_a^b\dfrac{\mathrm{d}x}{\ln x}<\dfrac{2b}{\ln b};
              \]
        \item 设 $0<a<b$,证明:
              \[
                  \left|\int_a^b\dfrac{\sin x}{x}\,\mathrm{d}x\right|\leq\dfrac{2}{a}.
              \]
    \end{enumerate}
\end{example}
\exampleblank

\begin{solution}
    \begin{enumerate}
        \item 当前题面不成立.令 $c=x$,则左端为
              \[
                  \left|\int_0^x\sin t^2\,\mathrm{d}t\right|.
              \]
              当 $x\to+\infty$ 时该式趋于非零的 Fresnel 积分值,而右端 $\dfrac{1}{x}\to0$.原题很可能应把上限 $x-c$ 写成 $x+c$,或补充其他限制.

        \item 当 $x>e^2$ 时,
              \[
                  \left(\dfrac{x}{\ln x}\right)'
                  =\dfrac{1}{\ln x}-\dfrac{1}{(\ln x)^2}
                  >\dfrac{1}{2\ln x}.
              \]
              因此
              \[
                  \int_a^b\dfrac{\mathrm{d}x}{\ln x}
                  <2\left[\dfrac{x}{\ln x}\right]_a^b
                  <\dfrac{2b}{\ln b}.
              \]

        \item 分部积分得
              \[
                  \begin{aligned}
                      \int_a^b\dfrac{\sin x}{x}\,\mathrm{d}x
                       & =\left[-\dfrac{\cos x}{x}\right]_a^b
                      -\int_a^b\dfrac{\cos x}{x^2}\,\mathrm{d}x.
                  \end{aligned}
              \]
              故
              \[
                  \left|\int_a^b\dfrac{\sin x}{x}\,\mathrm{d}x\right|
                  \leq\dfrac{1}{a}+\dfrac{1}{b}
                  +\int_a^b\dfrac{\mathrm{d}x}{x^2}
                  =\dfrac{2}{a}.
              \]
    \end{enumerate}
\end{solution}

%例题6.4.8
\begin{example}
    设 $f(x)$ 是 $\lbrack -\pi,\pi\rbrack$ 上的凸函数,$f'(x)$ 有界,求证:
    \[
        a_{2n}=\dfrac{1}{\pi}\int_{-\pi}^\pi f(x)\cos2nx\,\mathrm{d}x\geq0,
    \]
    \[
        a_{2n+1}=\dfrac{1}{\pi}\int_{-\pi}^\pi f(x)\cos(2n+1)x\,\mathrm{d}x\leq0.
    \]
\end{example}
\exampleblank

\begin{solution}
    对正整数 $k$,记
    \[
        I_k=\int_{-\pi}^{\pi}f(x)\cos kx\,\mathrm{d}x.
    \]
    凸函数的导数 $f'$ 单调不减且有界,因而可作 Riemann--Stieltjes 分部积分.先作普通分部积分,再对单调函数 $f'$ 分部积分,得到
    \[
        I_k=\dfrac{1}{k^2}
        \int_{-\pi}^{\pi}[(-1)^k-\cos kx]\,\mathrm{d}f'(x).
    \]
    测度 $\mathrm{d}f'\geq0$.当 $k=2n$ 时,括号内为 $1-\cos2nx\geq0$,故 $I_{2n}\geq0$;当 $k=2n+1$ 时,括号内为 $-1-\cos(2n+1)x\leq0$,故 $I_{2n+1}\leq0$.除以 $\pi$ 即得题设结论.
\end{solution}

%例题6.4.9
\begin{example}
    设 $f(x)\geq0$ 在 $\lbrack -\pi,\pi\rbrack$ 上单调递减,证明:
    \[
        a_{2n}=\dfrac{1}{\pi}\int_{-\pi}^\pi f(x)\sin2nx\,\mathrm{d}x\geq0,
    \]
    \[
        a_{2n+1}=\dfrac{1}{\pi}\int_{-\pi}^\pi f(x)\sin(2n+1)x\,\mathrm{d}x\leq0.
    \]
\end{example}
\exampleblank

\begin{solution}
    把正,负半轴配对.令
    \[
        h(x)=f(x)-f(-x),
        \qquad 0\leq x\leq\pi.
    \]
    由于 $f$ 单调递减,$h(0)=0$ 且 $h$ 单调递减.于是对正整数 $k$,
    \[
        \int_{-\pi}^{\pi}f(x)\sin kx\,\mathrm{d}x
        =\int_0^\pi h(x)\sin kx\,\mathrm{d}x.
    \]
    令正测度 $\mathrm{d}\mu=-\mathrm{d}h$.由 $h(x)=-\mu(\lbrack0,x\rbrack)$ 及换序积分,
    \[
        \int_0^\pi h(x)\sin kx\,\mathrm{d}x
        =\dfrac{1}{k}\int_0^\pi[(-1)^k-\cos kt]\,\mathrm{d}\mu(t).
    \]
    因此 $k=2n$ 时该积分非负,$k=2n+1$ 时非正.除以 $\pi$ 即得结论.
\end{solution}

%例题6.4.10
\begin{example}
    证明下列命题:
    \begin{enumerate}
        \item
              \[
                  \int_0^1\dfrac{\mathrm{d}x}{1+x^n}
                  =1-\dfrac{k^2}{n}+o(1),
                  \qquad n\to\infty;
              \]
        \item
              \[
                  \lim\limits_{x\to0}\dfrac{1}{x}\int_x^{2x}\sin\dfrac{1}{t}\,\mathrm{d}t=0;
              \]
        \item 当 $\alpha<2$ 时,
              \[
                  \lim\limits_{x\to0^+}\dfrac{1}{x^\alpha}\int_0^x\sin\dfrac{1}{t}\,\mathrm{d}t=0.
              \]
    \end{enumerate}
\end{example}
\exampleblank

\begin{solution}
    \begin{enumerate}
        \item 当前题面中的 $k^2$ 未定义.正确的一阶渐近式为
              \[
                  \int_0^1\dfrac{\mathrm{d}x}{1+x^n}
                  =1-\dfrac{\ln2}{n}+o\left(\dfrac{1}{n}\right).
              \]
              事实上,令 $x=e^{-\frac{t}{n}}$,则
              \[
                  \begin{aligned}
                      n\left(1-\int_0^1\dfrac{\mathrm{d}x}{1+x^n}\right)
                       & =\int_0^{+\infty}
                      \dfrac{e^{-\frac{t}{n}}}{e^t+1}\,\mathrm{d}t \\
                       & \longrightarrow\int_0^{+\infty}
                      \dfrac{\mathrm{d}t}{e^t+1}
                      =\ln2.
                  \end{aligned}
              \]
              因而原题中的 $k^2$ 应核对.

        \item 令 $t=xs$,则原式化为
              \[
                  \int_1^2\sin\dfrac{1}{xs}\,\mathrm{d}s.
              \]
              令 $\lambda=\dfrac{1}{x}$.由于
              \[
                  \dfrac{\mathrm{d}}{\mathrm{d}s}
                  \cos\dfrac{\lambda}{s}
                  =\dfrac{\lambda}{s^2}\sin\dfrac{\lambda}{s},
              \]
              分部积分得
              \[
                  \begin{aligned}
                      \int_1^2\sin\dfrac{\lambda}{s}\,\mathrm{d}s
                       & =\dfrac{1}{\lambda}
                      \left[s^2\cos\dfrac{\lambda}{s}\right]_1^2
                      -\dfrac{2}{\lambda}\int_1^2s\cos\dfrac{\lambda}{s}\,\mathrm{d}s,
                  \end{aligned}
              \]
              故其绝对值不超过 $\dfrac{8}{\lambda}=8x$,所求极限为 $0$.

        \item 令 $u=\dfrac{1}{t}$,则
              \[
                  \int_0^x\sin\dfrac{1}{t}\,\mathrm{d}t
                  =\int_{\frac{1}{x}}^{+\infty}\dfrac{\sin u}{u^2}\,\mathrm{d}u.
              \]
              分部积分给出
              \[
                  \left|\int_{\frac{1}{x}}^{+\infty}\dfrac{\sin u}{u^2}\,\mathrm{d}u\right|
                  \leq2x^2.
              \]
              因此当 $\alpha<2$ 时,
              \[
                  \left|\dfrac{1}{x^\alpha}
                  \int_0^x\sin\dfrac{1}{t}\,\mathrm{d}t\right|
                  \leq2x^{2-\alpha}\longrightarrow0.
              \]
    \end{enumerate}
\end{solution}
